 % 本文件由 scripts/assemble.py 自动装配，勿手改（改 parts/ 后重跑）
% 第9章 弯曲时空量子场论

% 第9章 Quantum Field Theory in Curved Spacetime
\chapter{弯曲时空量子场论}

\section{引言}

% ---- 书页 376 / p391 ----
没有人相信，就引力而言广义相对论就是最终的定论。奇点定理从理论内部提供了证据，表明这个理论在某种意义上是不完备的；然而更有说服力的是这样一个事实：广义相对论是一个经典理论，而世界在根本上是量子力学的。寻找一个能实际运作的量子引力理论，驱动着当今理论物理中的大量研究，这一路上我们也确实学到了许多，但令人信服的成功始终遥不可及。

广义相对论有两个组成部分：其一是时空曲率的框架及其对物质的影响，其二是度规响应能量-动量的动力学（由 Einstein 方程描述）。虽然还没有一个真正的量子引力理论，我们仍然可以取广义相对论的前一部分——物质场在弯曲时空背景上传播这一观念——并考虑那些物质场本身是量子力学的情况。换句话说，我们把度规取为固定的，而不让它服从什么动力学方程，然后研究这个弯曲时空中的量子场论（quantum field theory，QFT）。

弯曲时空量子场论研究中的划时代事件，是 Hawking 在 1976 年认识到黑洞其实并不真的黑，而是会发出热辐射，其温度是一个正比于表面引力（surface gravity）$\kappa$ 的\textbf{Hawking 温度}（Hawking temperature），
\begin{equation*}
\boxed{\;T = \frac{\kappa}{2\pi}.\;}\tag{9.1}
\end{equation*}
（回忆一下，我们的单位制取 $\hbar = c = k = 1$；Hawking 温度实际上正比于 $\hbar$，并反比于 Boltzmann 常数 $k$。）自这一非凡的发现以来，弯曲时空量子场论已被置于相当严格的理论基础之上，尽管人们普遍认为，它的适用范围离任何可能的实验探测都还相当遥远。Schwarzschild 黑洞的表面引力为 $\kappa = 1/4GM$，其 Hawking 温度可以写成
\begin{equation*}
T = \frac{1}{8\pi GM} = 1.2\times 10^{26}\,K\left(\frac{1\,\mathrm{g}}{M}\right) = 6.0\times 10^{-8}\,K\left(\frac{M_\odot}{M}\right),
\tag{9.2}
\end{equation*}
其中 $M_\odot \sim 10^{33}$ g 是太阳的质量。因此，来自一个现实的天体物理黑洞的辐射，其温度甚至比 3K 的宇宙微波背景还要低得多，因而基本上毫无被观测到的希望。

% ---- 书页 377 / p392 ----
然而，近来宇宙学方面的观测在某种程度上改变了这一局面。一个例子是看似发现了宇宙正在加速，这一结果最容易解释为非零真空能（vacuum energy）的证据（如第 8 章所讨论）。虽然真空能的大小仍然是一个深刻的谜，但有一点似乎是清楚的：理解量子力学的物质在弯曲时空中如何行为，将在任何对这一谜题的最终解决中扮演重要角色。另一个例子来自宇宙学扰动。对微波背景和大尺度结构的观测，为原初扰动的近无标度谱提供了强有力的证据，其中包括那些波长将远大于常规宇宙学中视界尺度的扰动。关于这些扰动起源的领先理论来自暴胀（inflation）。在暴胀场景中，宇宙学扰动起源于正在暴胀的宇宙中量子场的真空涨落。如果这幅图像是正确的，那么我们在 CMB 图中看到的，就是原初量子涨落被宇宙的膨胀大大拉伸之后留下的印记，而正是这些涨落经由引力不稳定性最终成长为今天我们所看到的星系和星系团。因此，至少可以说，宇宙学观测为弯曲时空量子场论的研究提供了强有力的动机。

即使没有这种经验上的动机，基于弯曲时空量子场论的思想实验，也已经在我们对量子引力的试探性探索中被证明是非常富有成果的。特别是，霍金辐射所预言的黑洞蒸发导致了信息丢失佯谬（information-loss paradox），我们将在下文讨论。由于直接触及量子引力问题的真实实验是如此困难，我们必须依靠聚焦于广义相对论与量子力学之间张力的思想实验，这很像 Einstein 当年试图调和经典动力学与电磁学的洛伦兹不变性时所用的思想实验。

带着这些考虑，本章的目标是对弯曲时空量子场论的一些观念和结果作一个简要的介绍。许多广义相对论入门书不讲这个主题，通常是因为学习广义相对论不应以熟悉平直时空中的普通量子场论为先决条件。然而令人愉快的事实是：熟悉平直时空中的量子场论，对于学习弯曲时空中的量子场论完全不是必需的。这是因为，在平直时空中最有趣、最有用的那些量子场论特征，与在弯曲时空中有趣、有用的那些特征几乎完全不同。往深处说，量子场论不过是量子力学系统的一个例子，就同方势阱或氦原子一样。一旦定义了一个场论，它在平直时空中的应用（对粒子物理或凝聚态物质）自然会聚焦于各种场之间相互作用的问题，通常把相互作用处理成围绕某个自然真空态的微扰。然而在弯曲时空中，我们感兴趣的通常是时空本身对场的影响，对此相互作用根本不是要点。因此我们可以考虑自由（无相互作用）场，但我们必须格外小心地定义什么才是合适的真空态。（的确，正如我们将要看到的，我们与之打交道的态几乎全都是真空态！）其结果是，平直时空量子场论的知识对眼下的讨论不仅不必要，甚至可能帮不上什么忙；唯一的先决条件是熟悉普通量子力学的基础。

% ---- 书页 378 / p393 ----
我们将循序渐进地走向弯曲时空中的量子场论，从复习这样一个系统的量子力学开始——每个物理学家在遇到棘手问题时都会求助于它：简谐振子（simple harmonic oscillator）。这当然是量子力学原理如何运作的一个典范性例子，而且还有一个额外的好处：当我们接下来转向场论时会发现，平直时空中自由场的量子力学恰好就是无穷多个谐振子的量子力学。（并不是空间每一点上都有一个振子，而是场的傅里叶变换中的每一个模式都表现得像一个谐振子。）于是，从谐振子到场论的过渡相当直截了当。一旦掌握了场论的基础，鉴于我们先前对广义相对论的学习，推广到弯曲时空并不十分困难，尽管沿途会遇到不少微妙之处。我们的讨论必然有些浅尝辄止，重心放在这样一个目标上：通过理解平直时空中的安鲁效应（Unruh effect）来理解霍金辐射的物理基础。特别是，我们不会讨论弯曲时空量子场论在宇宙学中的重要应用，也不会深入细致地考察重整化及相关问题。我们在很大程度上遵循 Birrell and Davies (1982) 的讨论；欲作进一步的探究，可查阅该书、Wald (1994)，或 Ford\footnote{L. H. Ford, “Quantum field theory in curved spacetime,” (1997), \texttt{http://arxiv.org/gr-qc/9707062}。} 的综述。

\section{量子力学}

量子场论只不过是量子力学系统的一个特定例子，所以我们可以从提醒自己这意味着什么开始。当然，尽管世界在根本上是量子力学的，我们的直觉却更容易同经典物理对上号，所以让我们先通过思考经典力学来搭好舞台。任何描述某个特定系统的物理理论，无论经典的还是量子的，都由以下三个问题的答案构成：

\begin{enumerate}
\item 系统的可能态是什么？在经典力学中，态空间通常由一组坐标和动量给出（我们可以把它们看作系统的“初始条件”）。它们可以被精确地指定，而这就是关于系统状态所要知道的一切。
\item 关于这个系统能观测到什么？这个问题在经典力学中往往只是隐含地处理，因为答案平淡无奇：坐标和动量的任何函数都有资格作为一个可观测量（observable）。
\item 系统如何演化？这通常由一组运动方程来表达。给定态和运动方程，随后的演化就唯一地确定了；于是，初始条件的空间就等价于该理论的经典解的空间。
\end{enumerate}

% ---- 书页 379 / p394 ----
为了让这些观念更具体，同时也因为它与我们研究场论直接相关，我们来考虑简谐振子。简谐振子可以想象成一维中受二次势作用的粒子。态由单个坐标 $x$ 和单个动量 $p$ 指定。为了得到运动方程，我们可以从拉格朗日量出发，它用 $x$ 及其时间导数 $\dot{x}$ 写成
\begin{equation*}
L = \frac{1}{2}\dot{x}^2 - \frac{1}{2}\omega^2 x^2,
\tag{9.3}
\end{equation*}
其中为方便起见，我们把振子的质量取为 1。我们可以立即导出运动方程
\begin{equation*}
\ddot{x} + \omega^2 x = 0.
\tag{9.4}
\end{equation*}
然而，为了过渡到量子力学，用哈密顿量来工作更加方便，它是 $x$ 和 $p$（而非 $x$ 和 $\dot{x}$）的函数。哈密顿量通过勒让德变换（Legendre transformation）与拉格朗日量相联系，
\begin{equation*}
H = p\dot{x} - L,
\tag{9.5}
\end{equation*}
其中动量满足
\begin{equation*}
p = \frac{\partial L}{\partial \dot{x}} = \dot{x}.
\tag{9.6}
\end{equation*}
于是我们得到振子的哈密顿量，
\begin{equation*}
H = \frac{1}{2}p^2 + \frac{1}{2}\omega^2 x^2,
\tag{9.7}
\end{equation*}
以及哈密顿方程
\begin{equation*}
\frac{dx}{dt} = \partial_p H = p, \qquad \frac{dp}{dt} = -\partial_x H = -\omega^2 x,
\tag{9.8}
\end{equation*}
它们充当运动方程。这些解当然是直截了当的；把它们表示成复数很有用：
\begin{equation*}
x(t) = x_0\, e^{i(\omega t + \alpha_0)},
\tag{9.9}
\end{equation*}
其中 $x_0$ 是振幅，$\alpha_0$ 是相位。到最后我们取实部，就得到物理答案。

现在我们转向量子力学。虽然量子力学与经典力学有着深刻的不同，但一个给定的理论仍然由与上面所列相同的三个问题的答案构成，只不过答案采取了多少有些不同的形式。

% ---- 书页 380 / p395 ----
\begin{enumerate}
\item 系统的态表示为\textbf{希尔伯特空间}（Hilbert space）中的一个元素。在数学上，希尔伯特空间就是一个装备了复数值内积的复矢量空间，该内积具有如下性质：以相反的次序取两个态的内积，等价于取复共轭。我们把希尔伯特空间的元素记作 $|\psi\rangle$，对偶空间的元素记作 $\langle\psi|$，于是 $|\psi_1\rangle$ 与 $|\psi_2\rangle$ 的内积就是 $\langle\psi_2|\psi_1\rangle$，并且满足
\begin{equation*}
\langle\psi_2|\psi_1\rangle^{*} = \langle\psi_1|\psi_2\rangle.
\tag{9.10}
\end{equation*}
（对于有关空间完备性的技术要求，我们在此不作深究。）在量子力学中，我们所感兴趣的希尔伯特空间往往都是无穷维的。例如，如果一个经典系统由坐标 $x$ 和动量 $p$ 来表示，那么希尔伯特空间可以取成由 $x$ 的所有平方可积复值函数组成，或者等价地由 $p$ 的所有平方可积复值函数组成（但不能同时取两者）。
\item 可观测量表示为希尔伯特空间上的\textbf{自伴算符}（self-adjoint operators）。“自伴”的定义实际上相当微妙，但在简单的情形下，它就相当于我们通常对厄米算符（Hermitian operator）的理解，
\begin{equation*}
A^{\dagger} = A,
\tag{9.11}
\end{equation*}
其中 $A^{\dagger}$ 满足
\begin{equation*}
\langle\psi_2|A\psi_1\rangle = \langle A^{\dagger}\psi_2|\psi_1\rangle
\tag{9.12}
\end{equation*}
对所有态 $|\psi_1\rangle$、$|\psi_2\rangle$ 成立。当然，许多算符并不是厄米的，但可观测量应当具有这一性质。一般说来，这类算符彼此不对易，因此我们无法同时精确地指明我们想要测量的关于系统的一切量的数值；将会存在一个对易可观测量完备集，它一次性地代表了我们对一个系统所能说的一切。
\item 系统的演化可以用两种方式之一来表示：或者作为希尔伯特空间中态矢量的幺正演化（\textbf{Schrödinger 绘景}，Schrödinger picture），或者保持态固定不动，而让可观测量按照运动方程演化（\textbf{Heisenberg 绘景}，Heisenberg picture）。
\end{enumerate}

严格说来，量子力学就是与经典力学不同；完全不必非得从一个经典模型出发，把它“量子化”。尽管如此，我们通常恰恰就是这么做的。即便是简单的经典模型，构造其量子化版本的方法也不止一种；其中包括正则量子化（canonical quantization）和路径积分量子化（path-integral quantization），以及一些更加奇特的手续。更糟的是，经典理论与量子理论之间并不存在简单的映射；有的经典理论没有良好定义的量子对应物，有的经典理论有好几个量子版本，还有的量子理论没有任何经典的类似物。就我们眼下的目的而言，我们可以对这些微妙之处满不在乎地一概置之不理，直接进行正则量子化。

% ---- 书页 381 / p396 ----
简谐振子再一次提供了一个有用的例子。先考虑大家熟悉的 Schrödinger 绘景，其中态由随时间演化的复值波函数表示，比如 $\psi(x,t)$。波函数实际上就是态矢量 $|\psi\rangle$ 的那一组分量，在“$\delta$ 函数位置基”$|x\rangle$ 中表出，于是 $|\psi(t)\rangle = \int dx\,\psi(x,t)|x\rangle$。正则量子化的内容就是加上正则对易关系（canonical commutation relation）
\begin{equation*}
[\hat{x}, \hat{p}] = i,
\tag{9.13}
\end{equation*}
把它加在坐标算符 $\hat{x}$ 及其共轭动量 $\hat{p}$ 上。对于由依赖 $x$ 和 $t$ 的波函数所表示的态，$\hat{x}$ 不过就是乘以 $x$，所以式 (9.13) 可以通过令
\begin{equation*}
\hat{p} = -i\partial_x.
\tag{9.14}
\end{equation*}
来实现。哈密顿算符是
\begin{equation*}
H = -\frac{1}{2}\partial_x^2 + \frac{1}{2}\omega^2 x^2,
\tag{9.15}
\end{equation*}
而运动方程就是 Schrödinger 方程，
\begin{equation*}
H\psi = i\partial_t\psi.
\tag{9.16}
\end{equation*}
由于哈密顿量不含时间，这个方程的解可以分离成空间的函数与时间的函数，$\psi(x,t) = f(t)g(x)$。这些解构成一个由整数 $n \geq 0$ 标记的离散集合，我们求得（在相差归一化因子的意义下）
\begin{equation*}
\psi_n(x,t) = e^{-(1/2)\omega x^2} H_n(\sqrt{\omega}\,x)\, e^{-iE_n t},
\tag{9.17}
\end{equation*}
其中 $H_n$ 是 $n$ 次厄米多项式（Hermite polynomial），而
\begin{equation*}
E_n = \left(n + \frac{1}{2}\right)\omega.
\tag{9.18}
\end{equation*}
这些态都是 $H$ 的本征函数，$E_n$ 就是能量本征值。振子的任意态无非就是能量本征态的叠加，
\begin{equation*}
\psi(x,t) = \sum_n c_n \psi_n(x,t),
\tag{9.19}
\end{equation*}
其中诸系数 $c_n$ 适当归一化。

这个简短的概述已包含了量子力学振子的许多重要特征。能量本征态具有分立谱；这正是它被称作“量子”力学的原因（尽管找出具有连续谱的系统并不困难）。存在一个能量最低的基态，以及一组由各自能量本征值唯一标记的激发态。基态具有非零的能量，
\begin{equation*}
E_0 = \frac{1}{2}\omega,
\tag{9.20}
\end{equation*}
它有时被称为“零点”能（zero-point energy）。有趣的是，经典系统的最小能量本来会是零，对应于 $x = 0$、$p = 0$ 的粒子。量子的零点能可以追溯到 Heisenberg 不确定性原理：它不许我们同时在位置和动量上把一个态定域化；因此振子中总存在一个最低限度的“抖动”，导致非零的基态能量。另一方面，我们当然也可以选择去考察一个势由 $V(x) = \frac{1}{2}\omega^2 x^2 - \frac{1}{2}\omega$ 给出的振子；我们的分析将完全相同，只不过式 (9.18) 中的因子 $\frac{1}{2}$ 不再出现，而基态能量将会是零。量子力学并不坚持非要有非零的零点能，它只是把能量从经典值那里挪动了一下。

% ---- 书页 382 / p397 ----
求解简谐振子的另一种办法，是引进产生算符和湮灭算符 $\hat{a}^{\dagger}$ 与 $\hat{a}$（常称为升算符和降算符，raising and lowering operators），其定义为
\begin{equation*}
\hat{a} = \frac{1}{\sqrt{2\omega}}\left(\omega\hat{x} + i\hat{p}\right), \qquad \hat{a}^{\dagger} = \frac{1}{\sqrt{2\omega}}\left(\omega\hat{x} - i\hat{p}\right),
\tag{9.21}
\end{equation*}
于是
\begin{equation*}
\hat{x} = \frac{1}{\sqrt{2\omega}}(\hat{a} + \hat{a}^{\dagger}), \qquad \hat{p} = -i\sqrt{\frac{\omega}{2}}(\hat{a} - \hat{a}^{\dagger}).
\tag{9.22}
\end{equation*}
利用我们先前关于对易关系的表达式 (9.13) 和哈密顿量 (9.7)，我们可以容易地算出产生算符和湮灭算符的对易关系，
\begin{equation*}
[\hat{a}, \hat{a}^{\dagger}] = 1,
\tag{9.23}
\end{equation*}
以及哈密顿量的新表达式，
\begin{equation*}
H = \left(\hat{a}^{\dagger}\hat{a} + \frac{1}{2}\right)\omega.
\tag{9.24}
\end{equation*}
产生/湮灭算符与哈密顿量按如下方式对易：
\begin{align*}
[H, \hat{a}] &= -\omega\hat{a} \\
[H, \hat{a}^{\dagger}] &= \omega\hat{a}^{\dagger}.
\tag{9.25}
\end{align*}
把这个版本的哈密顿量与能量本征值 (9.18) 相比较，我们受到启发，定义一个粒子数算符（number operator）
\begin{equation*}
\hat{n} = \hat{a}^{\dagger}\hat{a}.
\tag{9.26}
\end{equation*}

% ---- 书页 383 / p398 ----
让我们想一想，为什么产生/湮灭算符和粒子数算符配得上这些名字。考虑粒子数算符的一个本征态 $|n\rangle$，
\begin{equation*}
\hat{n}|n\rangle = n|n\rangle,
\tag{9.27}
\end{equation*}
其中左边的 $\hat{n}$ 代表粒子数算符，而右边第一个 $n$ 代表实际的数 $n$。（这个公式是整个量子力学中最迷人的公式。）摆弄一下对易关系，容易证明
\begin{align*}
\hat{n}\hat{a}^{\dagger}|n\rangle &= (n+1)\hat{a}^{\dagger}|n\rangle \\
\hat{n}\hat{a}|n\rangle &= (n-1)\hat{a}|n\rangle.
\tag{9.28}
\end{align*}
这样，当 $\hat{a}^{\dagger}$ 作用在 $|n\rangle$ 上时，它给出 $\hat{n}$ 的另一个本征态，本征值升高了 1；而 $\hat{a}$ 给出本征值降低了 1 的本征态。和前面一样，我们可以证明 $n$ 取从 $0$ 到 $\infty$ 的整数值，因此必定存在一个真空态 $|0\rangle$，满足
\begin{equation*}
\hat{a}|0\rangle = 0.
\tag{9.29}
\end{equation*}
从这个态出发，借助产生算符的逐次作用，我们可以构造出所有的本征态，
\begin{equation*}
|n\rangle = \frac{1}{\sqrt{n!}}\left(\hat{a}^{\dagger}\right)^{n}|0\rangle.
\tag{9.30}
\end{equation*}
粒子数算符计数的是高出基态的激发的数目。本征态的集合 $|n\rangle$ 充当一组基；任何态都是这些态的适当线性组合。产生算符和湮灭算符按如下方式作用在它们之上：
\begin{align*}
\hat{a}|n\rangle &= \sqrt{n}\,|n-1\rangle \\
\hat{a}^{\dagger}|n\rangle &= \sqrt{n+1}\,|n+1\rangle,
\tag{9.31}
\end{align*}
而每个态的能量当然由式 (9.18) 给出。基矢态取为不含时间，于是，一个服从 Schrödinger 方程的物理系统将用态
\begin{equation*}
|\psi(t)\rangle = \sum_n c_n e^{-iE_n t}|n\rangle,
\tag{9.32}
\end{equation*}
来描述，其中诸 $c_n$ 同样是常数系数。

为了让向场论的过渡更加平滑，把这种 Schrödinger 绘景的描述转译成 Heisenberg 绘景是有用的，在 Heisenberg 绘景中态是固定的，算符则随时间演化。给定 Schrödinger 方程 (9.16)，任何态都可以形式地写成某个固定的初始态经一个幺正时间演化算符作用后的结果，
\begin{equation*}
|\psi(t)\rangle = U(t)|\psi(0)\rangle,
\tag{9.33}
\end{equation*}

% ---- 书页 384 / p399 ----
其中
\begin{equation*}
U(t) = \mathcal{P}\, e^{-i\int H\,dt},
\tag{9.34}
\end{equation*}
（幺正的意思是 $U^{\dagger}U = 1$。）符号 $\mathcal{P}$ 代表路径排序（path-ordering），如附录 I 中所讨论。当然，如果哈密顿量不含时间，我们就简单地有 $U(t) = e^{-iHt}$。一个不含时算符 $A$ 在两个含时态 $|\psi_1(t)\rangle$ 和 $|\psi_2(t)\rangle$ 之间矩阵元的 Schrödinger 绘景表达式，于是可以写成 Heisenberg 绘景的表达式，即用含时算符 $A(t)$ 与不含时态表示为
\begin{align*}
\langle\psi_2(t)|A|\psi_1(t)\rangle &= \langle\psi_2(0)|U^{\dagger}(t)AU(t)|\psi_1(0)\rangle \\
&= \langle\psi_2|A(t)|\psi_1\rangle,
\tag{9.35}
\end{align*}
显然，其中 Heisenberg 绘景的算符由下式给出：
\begin{equation*}
A(t) = U^{\dagger}(t)AU(t).
\tag{9.36}
\end{equation*}
这样的算符满足\textbf{Heisenberg 运动方程}（Heisenberg equation of motion），
\begin{equation*}
\frac{dA(t)}{dt} = i[H, A(t)],
\tag{9.37}
\end{equation*}
它在这个绘景中取代了 Schrödinger 方程的位置。对于谐振子，我们会发现
\begin{equation*}
\frac{d\hat{a}}{dt} = -i\omega\hat{a}, \qquad \frac{d\hat{a}^{\dagger}}{dt} = i\omega\hat{a}^{\dagger},
\tag{9.38}
\end{equation*}
其解为
\begin{equation*}
\hat{a}(t) = e^{-i\omega t}\hat{a}(0), \qquad \hat{a}(t)^{\dagger} = e^{i\omega t}\hat{a}(0)^{\dagger}.
\tag{9.39}
\end{equation*}
由此我们立即得到
\begin{equation*}
\hat{n}(t) = \hat{a}(t)^{\dagger}\hat{a}(t) = \hat{a}(0)^{\dagger}\hat{a}(0),
\tag{9.40}
\end{equation*}
这反映了粒子数算符守恒这一事实。

人们常说，在 Heisenberg 绘景中态是不含时的；这话尽管不假，却有些令人困惑。或许更好的说法是：态延伸贯穿于整个时间，而不只是定义在某个固定时刻。为了把这一点看得更清楚，考虑一个受到外界影响的简谐振子，比如简单地在哈密顿量中加上一个强迫项，
\begin{equation*}
H = \frac{1}{2}p^2 + \frac{1}{2}\omega^2 x^2 + F(t),
\tag{9.41}
\end{equation*}
其中函数 $F(t)$ 在某个区间之外为零：
\begin{equation*}
F(t) =
\begin{cases}
0 & t < t_1 \\
F(t) & t_1 \leq t \leq t_2 \\
0 & t_2 < t .
\end{cases}
\tag{9.42}
\end{equation*}

% ---- 书页 385 / p400 ----
我们可以想象有人走过来，把我们的振子摇晃了一小会儿，然后就撒手不管了。在 Schrödinger 绘景中，我们会说：开始时处于基态的振子会被外力激发，末态不再是基态。而在 Heisenberg 绘景中，我们把态取成对所有时刻都成立的运动方程的解，并且说粒子数算符从零变成了某个别的值。

对于受到瞬态外力作用的振子，显然存在一组态，它们在早期看起来像能量本征态，尽管在将来就不再是这样了；我们可以把这样的态称为“in 态”$|n_{\mathrm{in}}\rangle$，它具有性质
\begin{equation*}
\hat{n}(t < t_1)|n_{\mathrm{in}}\rangle = n|n_{\mathrm{in}}\rangle.
\tag{9.43}
\end{equation*}
还有另外一组态，它们在晚期看起来像能量本征态，相应地称为“out 态”$|n_{\mathrm{out}}\rangle$，满足
\begin{equation*}
\hat{n}(t > t_2)|n_{\mathrm{out}}\rangle = n|n_{\mathrm{out}}\rangle.
\tag{9.44}
\end{equation*}
这两组态在所有时刻都存在，但只有在相应的渐近区域里，它们才看起来像能量本征态。任何一组都构成整个希尔伯特空间的一组基，所以特别地，我们可以用其中一组去展开另一组。例如，通过乘上一个 in 态的完备集，我们可以写出
\begin{equation*}
|n_{\mathrm{out}}\rangle = \sum_m \langle m_{\mathrm{in}}|n_{\mathrm{out}}\rangle|m_{\mathrm{in}}\rangle.
\tag{9.45}
\end{equation*}
复数 $\langle m_{\mathrm{in}}|n_{\mathrm{out}}\rangle$ 是矩阵元，原则上可以从哈密顿量 (9.41) 算出；它们合在一起构成\textbf{S 矩阵}（S-matrix）。一个装备了某种探测振子激发之手段的观者会发现，激发的数目被所施加的力改变了，而 S 矩阵编码了刻画渐近过去与渐近未来之间这些改变所需的全部信息。不用说，这整段讨论几乎可以原封不动地搬到场论中。在粒子物理里，外力的角色由不同粒子之间的相互作用扮演；而在我们的目的中，扮演这一角色的将是时空的曲率。

\section{平直时空中的量子场论}

正如我们已经提到的，量子场论只是量子力学系统的一个特定例子，其中被量子化的不是一个单个的振子，而是一个场（定义在时空上的一个函数，或者更一般地某个张量场）。我们从最简单的例子开始：平直时空中的自由标量场；从单个振子过渡到这个场论，只需再作寥寥几个推广。而把理论扩展到弯曲时空，照例是直截了当的：把理论写成协变形式，然后宣布它成立。然而，一旦失去了 Minkowski 空间的对称性，我们通常视为量子场论核心的一些观念就不再显得那么关键；特别是，“真空”与“粒子”的概念将失去它们的特权地位。（讲解量子力学的书偶尔会指出，波和粒子是适用范围各不相同的互补概念，但不要被误导；在量子场论中，真正基本的是场，而粒子只是在某些受限情形下才有用的近似概念。）本节我们先研究平直时空中的量子场论，下一节再推广到弯曲时空。

% ---- 书页 386 / p401 ----
我们从经典理论出发，具体说来是平直时空中的实标量场 $\phi(x^{\mu})$，正如我们在第 1 章中考虑过的那样，只是这次推广到 $n$ 维。作用量是拉格朗日密度（Lagrange density）的时空积分，$S = \int d^{\,n}x\,\mathcal{L}$；我们将考虑 Klein–Gordon 拉格朗日量
\begin{equation*}
\mathcal{L} = -\frac{1}{2}\eta^{\mu\nu}\partial_\mu\phi\,\partial_\nu\phi - \frac{1}{2}m^2\phi^2.
\tag{9.46}
\end{equation*}
没有必要纳入体积元因子 $\sqrt{|g|}$，因为我们在 Minkowski 空间中使用的是惯性坐标，其度规为
\begin{equation*}
ds^2 = -dt^2 + (d\mathbf{x})^2.
\tag{9.47}
\end{equation*}
运动方程是 Klein–Gordon 方程，
\begin{equation*}
\Box\phi - m^2\phi = 0.
\tag{9.48}
\end{equation*}
把场论转译成哈密顿描述是直截了当的。场的共轭动量（conjugate momentum）就是拉格朗日密度对该场的时间导数的导数，
\begin{equation*}
\pi = \frac{\partial\mathcal{L}}{\partial(\partial_0\phi)}.
\tag{9.49}
\end{equation*}
对于 Klein–Gordon 拉格朗日量 (9.46)，这就是
\begin{equation*}
\pi = \dot{\phi}.
\tag{9.50}
\end{equation*}
当然，一提到时间导数，就已默认我们选定了一个特定的惯性系；因此，哈密顿手续必然破坏显式的洛伦兹不变性。然而，只要我们足够小心，所得理论中的可观测量仍将是洛伦兹不变的。哈密顿量本身可以表示为哈密顿密度（Hamiltonian density）对空间的积分，
\begin{equation*}
H = \int d^{\,n-1}x\,\mathcal{H},
\tag{9.51}
\end{equation*}
哈密顿密度通过勒让德变换与拉格朗日密度相联系：
\begin{align*}
\mathcal{H}(\phi,\pi) &= \pi\dot{\phi} - \mathcal{L}(\phi,\partial_\mu\phi) \\
&= \frac{1}{2}\pi^2 + \frac{1}{2}(\nabla\phi)^2 + \frac{1}{2}m^2\phi^2,
\tag{9.52}
\end{align*}
其中 $(\nabla\phi)^2 = \delta^{ij}(\partial_i\phi)(\partial_j\phi)$。这个场论与谐振子之间的对应应当是清楚的：场值 $\phi(x)$ 扮演坐标 $x$ 的角色，而动量场 $\pi(x)$ 取代了单个的动量 $p$。态不再由某个固定时刻的两个数（$x$ 与 $p$）来规定，而是必须在某个固定时刻给出遍布空间的场值 $[\phi(x^i)$ 和 $\pi(x^i)]$ 作为初始数据，另外还多出一项振子情形所没有的梯度项；除此之外，这套形式体系是非常相似的。

% ---- 书页 387 / p402 ----
我们应当强调，$\phi(x^{\mu})$ \emph{不是}波函数；它是一个动力学变量，是谐振子情形中单个自由度 $x$ 的推广。如果对这个场论作 Schrödinger 绘景的量子化，我们会去定义一个复的波泛函（wave functional）$\Psi[\phi(x^{\mu})]$，用它表示发现场处于各个组态的概率振幅。然而，我们将改用 Heisenberg 绘景，于是我们首要关心的事情，就是把 $\phi$ 提升为一个量子算符。

首先，我们应当实际求出这个理论的解，以此完成经典分析。写下 Klein–Gordon 方程的解并不困难。一个好的例子是平面波，
\begin{equation*}
\phi(x^{\mu}) = \phi_0\, e^{ik_\mu x^\mu} = \phi_0\, e^{-i\omega t + i\mathbf{k}\cdot\mathbf{x}},
\tag{9.53}
\end{equation*}
其中波矢的分量是
\begin{equation*}
k^{\mu} = (\omega, \mathbf{k}),
\tag{9.54}
\end{equation*}
而频率必须满足色散关系（dispersion relation）
\begin{equation*}
\omega^2 = \mathbf{k}^2 + m^2.
\tag{9.55}
\end{equation*}
这样的解与式 (9.9) 给出的简谐振子的解之间存在着明显的相似性。但二者也有一个重要的差别：对于振子，只有一个独立的解。由于振子具有唯一的频率，当我们把两个具有给定振幅 $x_0$ 与相位 $\alpha_0$ 的解相加时，它们合并成第三个频率相同、但振幅和相位不同的解。在场论中这不再成立。给定式 (9.55)，频率就由空间波矢 $\mathbf{k}$ 决定，至少在相差正负号的意义下如此。因此，我们面对的不再是单一的一种解，而是由 $\mathbf{k}$ 与 $\omega$ 的正负号所参数化的一族解。

不过，我们仍然可以通过构造一组完备的、正交归一（orthonormal）的模式，来写下最一般的解，任何解都可以用这组模式表出。为了赋予“正交归一”以意义，我们需要在 Klein–Gordon 方程的解的空间上定义一个内积。虽然这些模式本身是时空的函数，但合适的内积可以表示成在常时超曲面 $\Sigma_t$ 上的积分：
\begin{equation*}
(\phi_1, \phi_2) = -i\int_{\Sigma_t}\left(\phi_1\partial_t\phi_2^{*} - \phi_2^{*}\partial_t\phi_1\right) d^{\,n-1}x.
\tag{9.56}
\end{equation*}

% ---- 书页 388 / p403 ----
正如我们所希望的，这个内积实际上\emph{不依赖于}积分所取的超曲面 $\Sigma_t$；利用 Stokes 定理和 Klein–Gordon 方程，你可以容易地验证这一点。把这个内积作用于两个波矢不同的平面波，得到
\begin{align*}
&\bigl(e^{ik_1^\mu x_\mu},\, e^{ik_2^\nu x_\nu}\bigr) \\
&\quad = -i\int_{\Sigma_t}\bigl(e^{-i\omega_1 t + i\mathbf{k}_1\cdot\mathbf{x}}\,\partial_t\, e^{i\omega_2 t - i\mathbf{k}_2\cdot\mathbf{x}} - e^{i\omega_2 t - i\mathbf{k}_2\cdot\mathbf{x}}\,\partial_t\, e^{-i\omega_1 t + i\mathbf{k}_1\cdot\mathbf{x}}\bigr)\, d^{\,n-1}x \\
&\quad = (\omega_2 + \omega_1)\, e^{-i(\omega_1 - \omega_2)t}\int_{\Sigma_t} e^{i(\mathbf{k}_1 - \mathbf{k}_2)\cdot\mathbf{x}}\, d^{\,n-1}x \\
&\quad = (\omega_2 + \omega_1)\, e^{-i(\omega_1 - \omega_2)t}\,(2\pi)^{n-1}\delta^{(n-1)}(\mathbf{k}_1 - \mathbf{k}_2),
\tag{9.57}
\end{align*}
其中我们用到了
\begin{equation*}
\int e^{i\mathbf{k}\cdot\mathbf{x}}\, d^{\,n-1}x = (2\pi)^{n-1}\delta^{(n-1)}(\mathbf{k}).
\tag{9.58}
\end{equation*}
于是，除非两个模式的空间波矢 $\mathbf{k}$——从而它们的频率 $\omega$——相等，内积就为零。因此，一组正交归一的模式解由下式给出：
\begin{equation*}
f_{\mathbf{k}}(x^{\mu}) = \frac{e^{ik_\mu x^\mu}}{\left[(2\pi)^{n-1}2\omega\right]^{1/2}},
\tag{9.59}
\end{equation*}
其中 $k^{\mu}$ 服从式 (9.55)，于是
\begin{equation*}
(f_{\mathbf{k}_1}, f_{\mathbf{k}_2}) = \delta^{(n-1)}(\mathbf{k}_1 - \mathbf{k}_2).
\tag{9.60}
\end{equation*}

给定色散关系 (9.55)，$\mathbf{k}$ 只能把频率确定到相差一个整体的正负号。我们的策略是坚持让 $\omega$ 永远为正数，并把复共轭 $f^{*}_{\mathbf{k}}(x^{\mu})$ 包括进来，以补全模式的集合。（复共轭会同时改变指数中 $\mathbf{k}$ 项与 $\omega$ 项的正负号，不过 $\mathbf{k}$ 的分量本来就定义在从 $-\infty$ 到 $\infty$ 的范围上。）$f_{\mathbf{k}}$ 模式被称为正频的（positive-frequency），意思是它们满足
\begin{equation*}
\partial_t f_{\mathbf{k}} = -i\omega f_{\mathbf{k}}, \qquad \omega > 0,
\tag{9.61}
\end{equation*}
而 $f^{*}_{\mathbf{k}}$ 模式是负频的（negative-frequency），满足
\begin{equation*}
\partial_t f^{*}_{\mathbf{k}} = i\omega f^{*}_{\mathbf{k}}, \qquad \omega > 0.
\tag{9.62}
\end{equation*}
（当心：这些模式虽然 $\omega > 0$，却被称作负频的，因为时间导数提出来的因子是 $+i\omega$ 而不是 $-i\omega$。）复共轭模式与原来的模式正交，
\begin{equation*}
(f_{\mathbf{k}_1}, f^{*}_{\mathbf{k}_2}) = 0,
\tag{9.63}
\end{equation*}
并彼此正交归一，但具有负范数，

\begin{equation*}
(f^{*}_{\mathbf{k}_1},\, f^{*}_{\mathbf{k}_2}) = -\delta^{(n-1)}(\mathbf{k}_1-\mathbf{k}_2).
\tag{9.64}
\end{equation*}

模式 $f_{\mathbf{k}}$ 与 $f^{*}_{\mathbf{k}}$ 合起来构成一个完备集，借助它我们可以展开 Klein--Gordon 方程的任何解。

要对这一理论进行正则量子化，我们把经典变量（场及其共轭动量）提升为作用在希尔伯特空间上的算符，并在等时超曲面上施加正则对易关系：

\begin{gather*}
[\phi(t,\mathbf{x}),\, \phi(t,\mathbf{x}')] = 0 \\
[\pi(t,\mathbf{x}),\, \pi(t,\mathbf{x}')] = 0 \\
[\phi(t,\mathbf{x}),\, \pi(t,\mathbf{x}')] = i\,\delta^{(n-1)}(\mathbf{x}-\mathbf{x}').
\tag{9.65}
\end{gather*}

在场论中，我们需要显式地声明场与其动量在整个空间中与自身对易；对单个谐振子而言这是隐含的，因为那里只有一个坐标和一个动量，各自必然与自身对易。$\delta$ 函数意味着等时的算符在空间中处处对易，只在空间点重合处例外；这一特征源于因果性的要求（类空间隔的算符不能相互影响）。

正如 Klein--Gordon 方程的经典解可以用模式 (9.59) 展开，量子算符场 $\phi(t,\mathbf{x})$ 也可以。把场算符模式展开的系数记作 $\hat{a}^{\dagger}_{\mathbf{k}}$ 和 $\hat{a}_{\mathbf{k}}$，我们有

\begin{equation*}
\phi(t,\mathbf{x}) = \int d^{\,n-1}k\,\bigl[\hat{a}_{\mathbf{k}}f_{\mathbf{k}}(t,\mathbf{x}) + \hat{a}^{\dagger}_{\mathbf{k}}f^{*}_{\mathbf{k}}(t,\mathbf{x})\bigr].
\tag{9.66}
\end{equation*}

把这一展开代入式 (9.65)，我们发现算符 $\hat{a}^{\dagger}_{\mathbf{k}}$ 和 $\hat{a}_{\mathbf{k}}$ 满足对易关系

\begin{gather*}
[\hat{a}_{\mathbf{k}},\, \hat{a}_{\mathbf{k}'}] = 0 \\
[\hat{a}^{\dagger}_{\mathbf{k}},\, \hat{a}^{\dagger}_{\mathbf{k}'}] = 0 \\
[\hat{a}_{\mathbf{k}},\, \hat{a}^{\dagger}_{\mathbf{k}'}] = \delta^{(n-1)}(\mathbf{k}-\mathbf{k}').
\tag{9.67}
\end{gather*}

于是这些算符满足产生和湮灭算符特有的对易关系，与简单谐振子的式 (9.23) 类似。当然，区别在于这样的算符有无穷多个，以 $\mathbf{k}$ 标记。由此可以看出把模式分为正频与负频的意义所在：正频模式是湮灭算符的系数，而负频模式是产生算符的系数。正、负频模式的概念以后会推广到静态时空，但不会推广到任意时空。

在谐振子的情形中，我们曾用产生和湮灭算符为希尔伯特空间定义了一组基，其中的基矢是粒子数算符（number operator）的本征态。同样的步骤也适用于自由标量场，只是现在我们必须对每个空间波矢 $\mathbf{k}$ 分别记录激发的数目。将会有一个唯一的真空态 $|0\rangle$，其特征是它能被每一个 $\hat{a}_{\mathbf{k}}$ 湮灭，

\begin{equation*}
\hat{a}_{\mathbf{k}}|0\rangle = 0 \qquad \text{for all } \mathbf{k}.
\tag{9.68}
\end{equation*}

具有 $n_{\mathbf{k}}$ 个相同动量 $\mathbf{k}$ 粒子的态，由 $\hat{a}^{\dagger}_{\mathbf{k}}$ 的反复作用产生，

\begin{equation*}
|n_{\mathbf{k}}\rangle = \frac{1}{\sqrt{n_{\mathbf{k}}!}}\left(\hat{a}^{\dagger}_{\mathbf{k}}\right)^{n_{\mathbf{k}}}|0\rangle,
\tag{9.69}
\end{equation*}

而具有各种动量 $\mathbf{k}_i$ 的 $n_i$ 个激发的态则为

\begin{equation*}
|n_1, n_2, \ldots, n_j\rangle = \frac{1}{\sqrt{n_1!n_2!\cdots n_j!}}\left(\hat{a}^{\dagger}_{\mathbf{k}_1}\right)^{n_1}\left(\hat{a}^{\dagger}_{\mathbf{k}_2}\right)^{n_2}\cdots\left(\hat{a}^{\dagger}_{\mathbf{k}_j}\right)^{n_j}|0\rangle.
\tag{9.70}
\end{equation*}

作用在这样的态上，产生和湮灭算符会如预期那样改变激发的数目：

\begin{align*}
\hat{a}_{\mathbf{k}_i}|n_1, n_2, \ldots, n_i, \ldots, n_j\rangle &= \sqrt{n_i}\,|n_1, n_2, \ldots, n_i-1, \ldots, n_j\rangle \\
\hat{a}^{\dagger}_{\mathbf{k}_i}|n_1, n_2, \ldots, n_i, \ldots, n_j\rangle &= \sqrt{n_i+1}\,|n_1, n_2, \ldots, n_i+1, \ldots, n_j\rangle.
\tag{9.71}
\end{align*}

我们可以对每个波矢定义一个粒子数算符，

\begin{equation*}
\hat{n}_{\mathbf{k}} = \hat{a}^{\dagger}_{\mathbf{k}}\hat{a}_{\mathbf{k}},
\tag{9.72}
\end{equation*}

它满足

\begin{equation*}
\hat{n}_{\mathbf{k}_i}|n_1, n_2, \ldots, n_i, \ldots, n_j\rangle = n_i|n_1, n_2, \ldots, n_i, \ldots, n_j\rangle.
\tag{9.73}
\end{equation*}

这些粒子数算符的本征态构成整个希尔伯特空间的一组基，称为 \textbf{Fock 基}（Fock basis）；由这组基构造的空间常被称作“Fock 空间”，但它当然就是原来的希尔伯特空间。

我们可能想考察的一件事，是我们的 Fock 基在洛伦兹变换下的行为。我们显然一直在利用 Minkowski 空间的对称性，例如用平面波作为 Klein--Gordon 方程解的一组基。这些模式的关键特征在于我们能够区分正频与负频，从而可以把 $\phi$ 的模式展开中的系数解释为湮灭和产生算符。现在考虑速度为 $\mathbf{v} = d\mathbf{x}/dt$ 的推动（boost），它给出新坐标 $x^{\mu'}$：

\begin{equation*}
t' = \gamma t - \gamma\,\mathbf{v}\cdot\mathbf{x}, \qquad \mathbf{x}' = \gamma\,\mathbf{x} - \gamma\,\mathbf{v}t,
\tag{9.74}
\end{equation*}

其中 $\gamma = 1/\sqrt{1-v^2}$，逆变换则为

\begin{equation*}
t = \gamma t' + \gamma\,\mathbf{v}\cdot\mathbf{x}', \qquad \mathbf{x} = \gamma\,\mathbf{x}' + \gamma\,\mathbf{v}t'.
\tag{9.75}
\end{equation*}

我们的模式函数在推动后的参考系中的时间导数为

\begin{align*}
\partial_{t'}f_{\mathbf{k}} &= \frac{\partial x^{\mu}}{\partial t'}\partial_{\mu}f_{\mathbf{k}} \\
&= \gamma(-i\omega)f_{\mathbf{k}} + \gamma\,\mathbf{v}\cdot(i\mathbf{k})f_{\mathbf{k}} \\
&= -i\omega' f_{\mathbf{k}}
\tag{9.76}
\end{align*}

其中

\begin{equation*}
\omega' = \gamma\omega - \gamma\,\mathbf{v}\cdot\mathbf{k}
\tag{9.77}
\end{equation*}

不过是在推动后的参考系中的频率。于是很明显，一个描述具有某些动量的粒子集合的态，被推动成描述同样一些粒子、但动量经过推动的态。因此，两个参考系中的总粒子数算符相一致，特别是真空态也相一致。在这个意义上，我们最初对惯性系的选择是无关紧要的。在下一节中我们将看到，我们之所以能找到正频和负频解，可以追溯到 Minkowski 时空中类时 Killing 矢量 $\partial_t$ 的存在；而 Fock 空间在基变换下的不变性，则可以追溯到所有这些类时 Killing 矢量都由洛伦兹变换相联系这一事实。因此，即使一个模式的频率依赖于惯性系的选择，正负频的分解却是不变的。

我们希望把哈密顿量

\begin{equation*}
H = \int d^{\,n-1}x\left[\frac{1}{2}\dot{\phi}^2 + \frac{1}{2}(\nabla\phi)^2 + \frac{1}{2}m^2\phi^2\right]
\tag{9.78}
\end{equation*}

用产生和湮灭算符表示出来，就像我们对谐振子做过的那样。我们可以逐项分析这个表达式，为简单起见从 $\phi^2$ 项开始：

\begin{align*}
&\frac{1}{2}m^2\int d^{\,n-1}x\,\phi^2 \\
&\quad = \frac{1}{2}m^2\int d^{\,n-1}x\, d^{\,n-1}k\, d^{\,n-1}k'\Bigl(\hat{a}_{\mathbf{k}}f_{\mathbf{k}} + \hat{a}^{\dagger}_{\mathbf{k}}f^{*}_{\mathbf{k}}\Bigr)\Bigl(\hat{a}_{\mathbf{k}'}f_{\mathbf{k}'} + \hat{a}^{\dagger}_{\mathbf{k}'}f^{*}_{\mathbf{k}'}\Bigr) \\
&\quad = \frac{1}{2}m^2\int d^{\,n-1}x\, d^{\,n-1}k\, d^{\,n-1}k'\Bigl(\hat{a}_{\mathbf{k}}\hat{a}_{\mathbf{k}'}f_{\mathbf{k}}f_{\mathbf{k}'} + \hat{a}^{\dagger}_{\mathbf{k}}\hat{a}_{\mathbf{k}'}f^{*}_{\mathbf{k}}f_{\mathbf{k}'} \\
&\qquad + \hat{a}_{\mathbf{k}}\hat{a}^{\dagger}_{\mathbf{k}'}f_{\mathbf{k}}f^{*}_{\mathbf{k}'} + \hat{a}^{\dagger}_{\mathbf{k}}\hat{a}^{\dagger}_{\mathbf{k}'}f^{*}_{\mathbf{k}}f^{*}_{\mathbf{k}'}\Bigr).
\tag{9.79}
\end{align*}

细看括号中的第一项，并暂时忽略对 $\mathbf{k}$ 的积分，代入模式函数 (9.59) 的显式形式，可得

\begin{align*}
\int d^{\,n-1}x\, d^{\,n-1}k'\,\hat{a}_{\mathbf{k}}\hat{a}_{\mathbf{k}'}f_{\mathbf{k}}f_{\mathbf{k}'}
&= \int d^{\,n-1}x\, d^{\,n-1}k'\,\hat{a}_{\mathbf{k}}\hat{a}_{\mathbf{k}'}\frac{e^{-i(\omega+\omega')t}e^{i(\mathbf{k}+\mathbf{k}')\cdot\mathbf{x}}}{2(2\pi)^{n-1}\sqrt{\omega\omega'}} \\
&= \int d^{\,n-1}k'\,\hat{a}_{\mathbf{k}}\hat{a}_{\mathbf{k}'}\frac{e^{-i(\omega+\omega')t}}{2\sqrt{\omega\omega'}}\,\delta^{(n-1)}(\mathbf{k}+\mathbf{k}') \\
&= \hat{a}_{\mathbf{k}}\hat{a}_{-\mathbf{k}}\frac{e^{-2i\omega t}}{2\omega},
\tag{9.80}
\end{align*}

这里我们再次用到了式 (9.58)。类似地计算式 (9.79) 中的其他项，我们发现对哈密顿量的势能贡献于是变为

\begin{align*}
\frac{1}{2}m^2\int d^{\,n-1}x\,\phi^2 = \frac{1}{2}m^2\int d^{\,n-1}k\left(\frac{1}{2\omega}\right)\Bigl[&\hat{a}_{\mathbf{k}}\hat{a}_{-\mathbf{k}}e^{-2i\omega t} \\
&+ \hat{a}^{\dagger}_{\mathbf{k}}\hat{a}_{\mathbf{k}} + \hat{a}_{\mathbf{k}}\hat{a}^{\dagger}_{\mathbf{k}} + \hat{a}^{\dagger}_{\mathbf{k}}\hat{a}^{\dagger}_{-\mathbf{k}}e^{2i\omega t}\Bigr].
\tag{9.81}
\end{align*}

对于动能和梯度能这两部分，求导会分别拉下 $\omega$ 和 $\mathbf{k}$ 的因子；我们得到

\begin{equation*}
\frac{1}{2}\int d^{\,n-1}x\,\dot{\phi}^2 = \frac{1}{2}\int d^{\,n-1}k\left(\frac{\omega}{2}\right)\Bigl[-\hat{a}_{\mathbf{k}}\hat{a}_{-\mathbf{k}}e^{-2i\omega t} + \hat{a}^{\dagger}_{\mathbf{k}}\hat{a}_{\mathbf{k}} + \hat{a}_{\mathbf{k}}\hat{a}^{\dagger}_{\mathbf{k}} - \hat{a}^{\dagger}_{\mathbf{k}}\hat{a}^{\dagger}_{-\mathbf{k}}e^{2i\omega t}\Bigr]
\tag{9.82}
\end{equation*}

以及

\begin{align*}
\frac{1}{2}\int d^{\,n-1}x\,(\nabla\phi)^2 = \frac{1}{2}\int d^{\,n-1}k\left(\frac{\mathbf{k}^2}{2\omega}\right)\Bigl[&\hat{a}_{\mathbf{k}}\hat{a}_{-\mathbf{k}}e^{-2i\omega t} \\
&+ \hat{a}^{\dagger}_{\mathbf{k}}\hat{a}_{\mathbf{k}} + \hat{a}_{\mathbf{k}}\hat{a}^{\dagger}_{\mathbf{k}} + \hat{a}^{\dagger}_{\mathbf{k}}\hat{a}^{\dagger}_{-\mathbf{k}}e^{2i\omega t}\Bigr].
\tag{9.83}
\end{align*}

利用 $\omega^2 = \mathbf{k}^2 + m^2$，我们可以把所有结果合在一起，把标量场理论的哈密顿量写成

\begin{align*}
H &= \frac{1}{2}\int d^{\,n-1}k\,\bigl[\hat{a}^{\dagger}_{\mathbf{k}}\hat{a}_{\mathbf{k}} + \hat{a}_{\mathbf{k}}\hat{a}^{\dagger}_{\mathbf{k}}\bigr]\omega \\
&= \int d^{\,n-1}k\,\Bigl[\hat{n}_{\mathbf{k}} + \frac{1}{2}\delta^{(n-1)}(0)\Bigr]\omega,
\tag{9.84}
\end{align*}

其中最后一步用到了对易关系 (9.67) 以及粒子数算符 $\hat{n}_{\mathbf{k}} = \hat{a}^{\dagger}_{\mathbf{k}}\hat{a}_{\mathbf{k}}$。按类似的逻辑，我们可以构造一个对应于总动量空间分量的算符，算出来是

\begin{equation*}
P^i = \int d^{\,n-1}k\,\hat{n}_{\mathbf{k}}k^i.
\tag{9.85}
\end{equation*}

正如我们所预期的，能量本征态将是具有固定激发数目的那些态，每个激发携带能量 $\omega$。Fock 基中的激发被解释为粒子。这就是粒子在量子场论中产生的方式：能量本征态是具有确定动量的粒子集合。当然，我们的模式是弥漫整个空间的平面波，而不是我们想到粒子时脑海中浮现的那种气泡室里的局域径迹。更糟的是，在弯曲时空中，波动方程将不再有我们能解释为粒子的、具有确定频率的平面波解。解决这两个问题的办法是操作性地思考，即从实验装置会观测到什么出发。最好的策略是先定义一个合理的粒子探测器（particle detector）概念，它在平直时空中还原为我们的直观图像，然后把“粒子”定义为“粒子探测器探测到的东西”。对一个定义恰当的粒子探测器，可以证明我们的平面波模式会像我们希望的那样“留下径迹”；在这样的探测器阵列中，如果一个平面波触发了一个探测器，那么它沿着从第一个探测器出发、由波矢给出的方向的路径，触发其他探测器的概率很高。（我们应该指出，如果你到 Fermilab 或 CERN 这样的地方参观一座真正的粒子加速器，会发现那些探测器与研究弯曲时空量子场论的理论家们所发明的探测器几乎毫无相似之处；不过究其本质，二者存在根本的相似性。）关于粒子探测器的讨论可参见 Birrell 和 Davies (1982)。

你或许会为哈密顿量 (9.84) 中的因子 $\delta^{(n-1)}(0)$ 感到担心——而且你理应如此。它意味着即使在真空态 $|0\rangle$ 下测量，哈密顿量也是无穷大。这一项是谐振子零点能 (9.20) 在场论中的对应物。在第 4 章讨论宇宙学常数时我们提到过，量子涨落会引起经典真空能形式上无穷大的位移；当把引力包括进来时，标量场哈密顿量中的这个无穷大贡献将表现为一个发散的宇宙学常数。它是无穷大量 $\delta^{(n-1)}(0)$ 在无穷大的 $\mathbf{k}$ 范围上的积分这一事实，可以转述为：总能量是对无限大空间中无穷大能量密度的积分。但真正的问题在于高频模式贡献的能量密度，而不是无限大的体积；如果我们把计算放进体积为 $L^{n-1}$ 的盒子里来做正则化，就会发现

\begin{equation*}
\frac{1}{2}\int d^{\,n-1}k\,\delta^{(n-1)}(0)\,\omega \to \frac{1}{2}\left(\frac{L}{2\pi}\right)^{n-1}\sum_{\mathbf{k}}\omega,
\tag{9.86}
\end{equation*}

即使 $L$ 有限它也发散，因为 $\mathbf{k}$（从而 $\omega$）可以任意大。在某个高动量 $k_{\mathrm{max}}$ 处放一个截断，就会回到式 (4.104)。

在简单谐振子的情形中我们指出过，当初要是选一个具有负极小值的经典势，零点能本可以避免；量子力学的贡献不一定代表真正的答案，只代表能量相对其经典值的位移。场论中也是如此；我们可以自由地定义最初的经典标量场理论，使得量子力学真空能为零。然而，我们不能逐模式地减去有限的能量，因为我们的自由仅仅是往势里、从而往哈密顿量密度 (9.52) 中加一个单独的常数。要得到真空态的有限哈密顿量，这个常数必须是无穷大。减去一个无穷大常数并没有什么不对；这是量子场论中一种历史悠久的技术，称为“重整化”（renormalization）。重整化有时看上去吓人，或者多少有些不正当，但实际上它完全合理；无穷大只出现在量子理论与它们经典对应物之间的关系之中，而不出现在任何可观测量之中。既然大自然多半既不知道也不在乎我们对经典力学的偏爱，重整化就不应该有什么特别令人不安的地方。

当然，一旦我们重整化得到有限的真空能，这个能量就可以取我们喜欢的任何值；它完全任意。这一点在弯曲时空量子场论中同样成立；我们也许无法把场分解成具有确定频率的模式，因而无法给每个模式指派一个真空能贡献，但仔细的分析允许我们把真空能重整化成你喜欢的任何数。同样，并没有发生什么深刻的事情；真空能在我们的经典模型中本来就是完全任意的，我们只是为了方便才把它选成零。宇宙学常数问题并不是因为量子力学贡献了大量真空能而产生的，因为这个贡献可以直接重整化掉；问题在于，所得到的任意数没有理由接近零。如前所述，从有效场论的角度看，问题还要更尖锐一些，因为对真空能的尺度存在一个逻辑上的预期，那就是未知的量子引力效应应当开始起作用的 Planck 标度。不过，贯穿本章我们只关心量子场在固定时空背景中的传播，而不把量子能动张量用作 Einstein 方程的源；因此我们可以选择无视宇宙学常数问题。

\section{弯曲时空中的量子场论}

在第 4 章中我们讨论过，把物理理论从平直时空推广到弯曲时空是多么容易——我们只需把理论写成坐标不变的形式，然后断言当时空弯曲时它们仍然成立。这一步骤对量子场论依然有效，尽管我们将不得不放弃一些在平直时空中显得不可或缺的概念。

我们从弯曲时空中标量场的拉格朗日密度出发，

\begin{equation*}
\mathcal{L} = \sqrt{-g}\left(-\frac{1}{2}g^{\mu\nu}\nabla_{\mu}\phi\nabla_{\nu}\phi - \frac{1}{2}m^2\phi^2 - \xi R\phi^2\right).
\tag{9.87}
\end{equation*}

除了可以预料到的度规 $g_{\mu\nu}$ 及其行列式的出现之外，我们还包含了与标量曲率 $R$ 的直接耦合，由一个常数 $\xi$ 参数化。由于 $\xi$ 无量纲，没有理由预期它很小；实际上，它自然应该是 1 的量级。文献中对 $\xi$ 的值有两种受青睐的选择：\textbf{最小耦合}（minimal coupling）干脆关掉与 $R$ 的直接相互作用，

\begin{equation*}
\xi = 0,
\tag{9.88}
\end{equation*}

而\textbf{共形耦合}（conformal coupling）则取

\begin{equation*}
\xi = \frac{(n-2)}{4(n-1)},
\tag{9.89}
\end{equation*}

在四维时它就是 $\xi = \frac{1}{6}$。利用附录 G 中的公式容易验证，当 $\xi$ 取这个值且 $m=0$ 时，标量场理论在共形变换 $g_{\mu\nu}\to\omega^2(x)g_{\mu\nu}$ 下不变。实际上，在真实世界中并没有充分的理由在最小耦合和共形耦合之间择取其一；最小耦合不增强任何对称性，而共形不变性显然不是大多数物理理论所具有的对称性。（由于共形变换是局部的标度改变，以质量这样的有量纲参数为特征的理论一般不会共形不变。）即使经典理论是共形不变的，量子化也可能破坏这一对称性，例如在与无质量夸克耦合的量子色动力学（QCD）理论中就发生了这种情况。一般而言，在四维中很难找到严格共形不变的相互作用理论，尽管已知某些具有高度超对称性的模型是共形不变的。

我们可以像先前那样继续量子化这个理论。共轭动量是

\begin{equation*}
\pi = \frac{\partial\mathcal{L}}{\partial(\nabla_0\phi)},
\tag{9.90}
\end{equation*}

对于拉格朗日量 (9.87)，它为

\begin{equation*}
\pi = \sqrt{-g}\,\nabla_0\phi.
\tag{9.91}
\end{equation*}

我们可以施加正则对易关系

\begin{gather*}
[\phi(t,\mathbf{x}),\, \phi(t,\mathbf{x}')] = 0 \\
[\pi(t,\mathbf{x}),\, \pi(t,\mathbf{x}')] = 0 \\
[\phi(t,\mathbf{x}),\, \pi(t,\mathbf{x}')] = \frac{i}{\sqrt{-g}}\,\delta^{(n-1)}(\mathbf{x}-\mathbf{x}').
\tag{9.92}
\end{gather*}

标量场的运动方程是

\begin{equation*}
\Box\phi - m^2\phi - \xi R\phi = 0.
\tag{9.93}
\end{equation*}

对具有诱导度规 $\gamma_{ij}$ 和单位法矢量 $n^{\mu}$ 的类空超曲面 $\Sigma$，该方程解上的内积为

\begin{equation*}
(\phi_1, \phi_2) = -i\int_{\Sigma}\left(\phi_1\nabla_{\mu}\phi^{*}_2 - \phi^{*}_2\nabla_{\mu}\phi_1\right)n^{\mu}\sqrt{\gamma}\, d^{\,n-1}x,
\tag{9.94}
\end{equation*}

它与 $\Sigma$ 的选取无关。

到目前为止一切顺利。为了继续我们在平直空间中所采取的步骤，现在应当引进一组正频和负频模式，它们构成式 (9.93) 解的完备基，把场算符 $\phi$ 用这些模式展开，并把算符系数解释为产生和湮灭算符。正是在这一步，我们的程序失灵了。由于一般不存在任何类时 Killing 矢量，我们通常无法找到能分离成时间依赖和空间依赖因子的波动方程的解，相应地也就不能把模式分类为正频或负频。我们可以找到一组基模式，但问题在于这样的模式组一般有很多，没有任何办法偏爱其中一组而舍弃其他，而真空和粒子数算符的概念将敏感地依赖于我们选择哪一组。

看看我们能做些什么。我们总能找到式 (9.93) 的一组正交归一的解 $f_i(x^{\mu})$，

\begin{equation*}
(f_i, f_j) = \delta_{ij},
\tag{9.95}
\end{equation*}

以及具有负范数的相应共轭模式，

\begin{equation*}
(f^{*}_i, f^{*}_j) = -\delta_{ij}.
\tag{9.96}
\end{equation*}

指标 $i$ 可以连续或分立；目前我们采用适合分立情形的记号。这些模式可以选成完备集，于是我们可以把场展开为

\begin{equation*}
\phi = \sum_i\left(\hat{a}_i f_i + \hat{a}^{\dagger}_i f^{*}_i\right).
\tag{9.97}
\end{equation*}

系数 $\hat{a}_i$ 和 $\hat{a}^{\dagger}_i$ 满足对易关系

\begin{gather*}
[\hat{a}_i,\, \hat{a}_j] = 0 \\
[\hat{a}^{\dagger}_i,\, \hat{a}^{\dagger}_j] = 0 \\
[\hat{a}_i,\, \hat{a}^{\dagger}_j] = \delta_{ij}.
\tag{9.98}
\end{gather*}

将存在一个真空态 $|0_f\rangle$，它被所有湮灭算符湮灭，

\begin{equation*}
\hat{a}_i|0_f\rangle = 0 \qquad \text{for all } i.
\tag{9.99}
\end{equation*}

从这个真空态出发，我们可以为希尔伯特空间定义整套 Fock 基。与之前一样，具有 $n_i$ 个激发的态由 $\hat{a}^{\dagger}_i$ 的反复作用产生，

\begin{equation*}
|n_i\rangle = \frac{1}{\sqrt{n_i!}}\left(\hat{a}^{\dagger}_i\right)^{n_i}|0_f\rangle,
\tag{9.100}
\end{equation*}

对具有不同种类激发的态也同样如此。我们甚至可以对每个模式定义一个粒子数算符，

\begin{equation*}
\hat{n}_{fi} = \hat{a}^{\dagger}_i\hat{a}_i.
\tag{9.101}
\end{equation*}

真空态和粒子数算符下标中的 $f$ 提醒我们，它们是相对于模式组 $f_i$ 定义的。

这套架构看上去与我们在平直空间中的非常相似；为什么我们不能宣布由 $\hat{a}^{\dagger}_i$ 产生的激发就是粒子，然后就此收工呢？我们本可以这样做，但必须面对这样一个事实：还存在其他可能的选择；基模式 $f_i(x^{\mu})$ 高度不唯一。考虑另外一组模式 $g_i(x^{\mu})$，它具有我们原来的模式所具有的全部性质，包括（连同共轭模式 $g^{*}_i$）构成一个完备基，相对于它可以展开我们的场算符，

\begin{equation*}
\phi = \sum_i\left(\hat{b}_i g_i + \hat{b}^{\dagger}_i g^{*}_i\right).
\tag{9.102}
\end{equation*}

湮灭和产生算符 $\hat{b}_i$ 与 $\hat{b}^{\dagger}_i$ 满足对易关系

\begin{gather*}
[\hat{b}_i,\, \hat{b}_j] = 0 \\
[\hat{b}^{\dagger}_i,\, \hat{b}^{\dagger}_j] = 0 \\
[\hat{b}_i,\, \hat{b}^{\dagger}_j] = \delta_{ij},
\tag{9.103}
\end{gather*}

并且将存在一个被所有湮灭算符湮灭的真空态 $|0_g\rangle$，

\begin{equation*}
\hat{b}_i|0_g\rangle = 0 \qquad \text{for all } i.
\tag{9.104}
\end{equation*}

我们可以通过在这个真空上反复作用产生算符来构造 Fock 基，并定义粒子数算符

\begin{equation*}
\hat{n}_{gi} = \hat{b}^{\dagger}_i\hat{b}_i.
\tag{9.105}
\end{equation*}

在从平直时空到弯曲时空的转变中，我们失去的是偏爱某一组模式甚于其他任何一组的任何理由。在平直时空中，我们能够通过要求模式相对于时间坐标为正频〔如式 (9.61) 所定义的〕而挑出一组自然的模式。时间坐标并不唯一，因为我们可以自由地做洛伦兹变换；但我们看到过，真空态和总粒子数算符在这样的变换下不变。因此，每个惯性系观者对什么是真空态、周围有多少粒子，都会有一致的看法。

在我们现在考虑的更一般情形中，如果一个观者相对于模式组 $f_i$ 定义粒子，而另一个观者使用模式组 $g_i$，那么他们对观测到多少个粒子（甚至究竟有没有观测到粒子）通常会有分歧。为了看清这一点，方便的做法是把每组模式用另一组展开，

\begin{align*}
g_i &= \sum_j\left(\alpha_{ij}f_j + \beta_{ij}f^{*}_j\right) \\
f_i &= \sum_j\left(\alpha^{*}_{ji}g_j - \beta_{ji}g^{*}_j\right).
\tag{9.106}
\end{align*}

从一组基模式到另一组的变换称为 \textbf{Bogoliubov 变换}（Bogoliubov transformation），实现该变换的矩阵 $\alpha_{ij}$ 和 $\beta_{ij}$ 称为 Bogoliubov 系数。利用模式函数的正交归一性，它们可以表示为

\begin{gather*}
\alpha_{ij} = (g_i, f_j) \\
\beta_{ij} = -(g_i, f^{*}_j).
\tag{9.107}
\end{gather*}

它们满足自己的归一化条件，

\begin{align*}
\sum_j\left(\alpha_{ik}\alpha^{*}_{jk} - \beta_{ik}\beta^{*}_{jk}\right) &= \delta_{ij} \\
\sum_j\left(\alpha_{ik}\beta_{jk} - \beta_{ik}\alpha_{jk}\right) &= 0.
\tag{9.108}
\end{align*}

Bogoliubov 系数除了描述模式之间的变换，还可以用来在算符之间变换，

\begin{align*}
\hat{a}_i &= \sum_j\left(\alpha_{ji}\hat{b}_j + \beta^{*}_{ji}\hat{b}^{\dagger}_j\right) \\
\hat{b}_i &= \sum_j\left(\alpha^{*}_{ij}\hat{a}_j - \beta^{*}_{ij}\hat{a}^{\dagger}_j\right).
\tag{9.109}
\end{align*}

现在设想系统处于 $f$ 真空 $|0_f\rangle$ 中，其中不会观测到任何 $f$ 粒子；我们想知道使用 $g$ 模式的观者会观测到多少粒子。因此我们来计算 $g$ 粒子数算符在 $f$ 真空中的期望值：

\begin{align*}
\langle 0_f|\hat{n}_{gi}|0_f\rangle &= \langle 0_f|\hat{b}^{\dagger}_i\hat{b}_i|0_f\rangle \\
&= \left\langle 0_f\left|\sum_{jk}\left(\alpha_{ij}\hat{a}^{\dagger}_j - \beta_{ij}\hat{a}_j\right)\left(\alpha^{*}_{ik}\hat{a}_k - \beta^{*}_{ik}\hat{a}^{\dagger}_k\right)\right|0_f\right\rangle \\
&= \sum_{jk}\left(-\beta_{ij}\right)\left(-\beta^{*}_{ik}\right)\langle 0_f|\hat{a}_j\hat{a}^{\dagger}_k|0_f\rangle \\
&= \sum_{jk}\beta_{ij}\beta^{*}_{ik}\langle 0_f|\left(\hat{a}^{\dagger}_k\hat{a}_j + \delta_{jk}\right)|0_f\rangle \\
&= \sum_{jk}\beta_{ij}\beta^{*}_{ik}\delta_{jk}\langle 0_f|0_f\rangle \\
&= \sum_j\beta_{ij}\beta^{*}_{ij}.
\tag{9.110}
\end{align*}

于是，$f$ 真空中 $g$ 粒子的数目可以用 Bogoliubov 系数表示为

\begin{equation*}
\langle 0_f|\hat{n}_{gi}|0_f\rangle = \sum_j|\beta_{ij}|^2.
\tag{9.111}
\end{equation*}

没有理由认为它为零；从一个视角看是空无一物的真空，从另一个视角看却翻腾着粒子。只要任何一个 $\beta_{ij}$ 非零，两个真空态就不会重合。为什么会这样，看看式 (9.109) 就能明白：$\beta_{ij}$ 描述的是一个基中的产生算符混入另一个基中湮灭算符的成分。

这些关于模式和粒子数算符的讨论可能显得不必要地抽象；诚然，如果一个真正的粒子探测器在可能是弯曲的时空中沿某条轨迹运动，它要么探测到粒子，要么探测不到，并不知道我们的场论用的是哪一组基模式。我们怎么知道这样的探测器实际使用的是哪一种“粒子”定义？答案在于：探测器测量的是沿其轨迹的固有时 $\tau$，并且将相对于该固有时来定义正频和负频。因此，如果能找到一组满足

\begin{equation*}
\frac{D}{d\tau}f_i = -i\omega f_i,
\tag{9.112}
\end{equation*}

的模式 $f_i$，我们就能用这些模式计算探测器将看到多少粒子。当然，一般不可能在整个时空中都找到这样的模式。有可能做到的一种情形是\emph{静态}（static）时空，此时我们有一个超曲面正交的类时 Killing 矢量 $K^{\mu}$。在那种情况下，我们可以选取这样的坐标：度规分量与时间坐标 $t$ 无关，并且不存在时间-空间交叉项：

\begin{equation*}
\partial_0 g_{\mu\nu} = 0, \qquad g_{0i} = 0.
\tag{9.113}
\end{equation*}

（指标 $i$、$j$ 现在指空间分量，不是模式标记。）对这样的度规，d'Alembert 算符作用在某个模式函数 $f(t,\mathbf{x})$ 上的结果是

\begin{equation*}
\Box f = \left[g^{00}\partial_0^2 + \frac{1}{2}g^{00}g^{ij}(\partial_i g_{00})\partial_j + g^{ij}\partial_i\partial_j - g^{ij}\Gamma^{k}_{ij}\partial_k\right]f.
\tag{9.114}
\end{equation*}

于是运动方程 (9.93) 可以写成

\begin{equation*}
\partial_0^2 f = -\left(g^{00}\right)^{-1}\left[g^{ij}\partial_i\partial_j + \frac{1}{2}g^{00}g^{ij}(\partial_i g_{00})\partial_j - g^{ij}\Gamma^{k}_{ij}\partial_k - \left(m^2 + \xi R\right)\right]f.
\tag{9.115}
\end{equation*}

左边的算符是纯时间导数，而右边的算符只涉及空间导数和仅依赖空间的函数。因此我们可以找到分离变量的解

\begin{equation*}
f_{\omega}(t,\mathbf{x}) = e^{-i\omega t}\bar{f}_{\omega}(\mathbf{x}),
\tag{9.116}
\end{equation*}

它们可以称为正频的，

\begin{equation*}
\partial_t f_{\omega}(t,\mathbf{x}) = -i\omega f_{\omega}(t,\mathbf{x}), \qquad \omega > 0.
\tag{9.117}
\end{equation*}

这个关系可以改写成坐标不变的形式

\begin{equation*}
\mathcal{L}_K f_{\omega} = K^{\mu}\partial_{\mu}f_{\omega} = -i\omega f_{\omega}, \qquad \omega > 0,
\tag{9.118}
\end{equation*}

其中 $\mathcal{L}_K f_{\omega}$ 表示 $f_{\omega}$ 沿 $K$ 的李导数。同时也存在负频的共轭模式，

\begin{equation*}
\mathcal{L}_K f^{*}_{\omega} = K^{\mu}\partial_{\mu}f^{*}_{\omega} = i\omega f^{*}_{\omega}, \qquad \omega > 0.
\tag{9.119}
\end{equation*}

模式 $(f_{\omega}, f^{*}_{\omega})$ 合起来将构成静态背景中波动方程解的一组基。除非这些模式与我们的探测器相关，否则它们的存在帮不了我们；如果探测器的轨迹沿着 Killing 场的轨道行进（四速度 $U^{\mu} = dx^{\mu}/d\tau$ 与 $K^{\mu}$ 成正比），固有时就将与 Killing 时间 $t$ 成正比，而相对于这个 Killing 矢量为正频的模式便可以作为描述 Fock 空间的自然基。在下一节讨论安鲁效应时，我们将看到这一现象的实际运作。

在上一节中我们提到过，量子场论中需要对真空能进行重整化。这一需求在弯曲时空中依然存在，但由于不存在偏好的模式基，构造恰当的重整化程序更加困难。不过，至少在某些情形下，人们已经发展出代数方法来严格地定义重整化的能动张量；我们不会深入细节地探讨这个专题，但至少应当介绍一下其背后的基本思想。基本想法是：即使存在曲率，在足够小的尺度上时空看上去应当是 Minkowski 式的。因为我们在平直时空中发现的真空能发散源于短波长模式，所以我们应当能把弯曲时空中场在非常小尺度上的行为与平直时空中的行为相匹配，并减去出现的任何发散。特别地，我们考虑量子场 $\phi$ 在某个态 $|\psi\rangle$ 中的两点函数，
\begin{equation*}
G(x_1, x_2) = \langle\psi|\phi(x_1)\phi(x_2)|\psi\rangle,
\tag{9.120}
\end{equation*}

其中 $x_1$ 与 $x_2$ 是两个时空点。当 $x_1$ 与 $x_2$ 彼此靠近时，Minkowski 真空中的两点函数会变得奇异。我们希望刻画这种奇异性，并坚持要求它对弯曲时空中的任何正则（regular）态都成立。所谓“彼此靠近”，指的是连接两点的最短测地线距离的平方 $\sigma(x_1,x_2)$ 趋于零。在 $x_1$ 与 $x_2$ 非常接近的极限下，测地线距离的平方就是

\begin{equation*}
\sigma(x_1, x_2) = g_{\mu\nu}(x_1^{\mu} - x_2^{\mu})(x_1^{\nu} - x_2^{\nu}), \qquad x_1 \to x_2.
\tag{9.121}
\end{equation*}

当然，在洛伦兹流形上，测地线距离不仅在两点重合时为零，在两点类光分离时也为零。因此我们引入一个小虚部，并取它趋于零的极限，办法是定义

\begin{equation*}
\sigma_{\epsilon}(x_1, x_2) = \sigma(x_1, x_2) + 2i\epsilon(t_1 - t_2) + \epsilon^{2}.
\tag{9.122}
\end{equation*}

这里 $t$ 是类时坐标，并且默认取 $\epsilon \to 0^{+}$ 的极限。（这个公式显而易见的坐标依赖性在此极限下无关紧要。）结果表明，Minkowski 时空中的自然真空具有一种唯一的奇异性结构：两点函数（在四维情形）包含一个形如 $1/(4\pi^{2}\sigma_{\epsilon})$ 的领头奇异性，以及一个正比于 $\ln\sigma_{\epsilon}$ 的次领头奇异性，其余各项都是正则的。因此我们要求弯曲时空中任何物理上合理的量子态都满足

\begin{equation*}
G(x_1, x_2) = \frac{U(x_1, x_2)}{4\pi^{2}\sigma_{\epsilon}} + V(x_1, x_2)\ln\sigma_{\epsilon} + W(x_1, x_2),
\tag{9.123}
\end{equation*}

其中函数 $U(x_1,x_2)$、$V(x_1,x_2)$ 和 $W(x_1,x_2)$ 在 $x_1=x_2$ 处都是正则的，并且 $U(x,x)=1$。具有这种性质的态称为\textbf{Hadamard 态}（Hadamard state）。可以证明，重整化的能动张量在所有 Hadamard 态中都有良好定义且非奇异；进一步还可以证明，它在任何非 Hadamard 态中都会奇异。如果 Hadamard 条件在某个部分 Cauchy 面上成立，那么它在依赖域内也将处处成立；换言之，能动张量可能在某个视界上变得奇异，但不会在某种适定初值数据的 Cauchy 发展（Cauchy development）内部变得奇异。因此，这种形式的态看来适合作为弯曲时空量子场论考虑的对象。详见 Wald (1994)。

我们看到，弯曲时空中的量子场论具备平直时空量子场论的大部分基本特征；关键的区别在于我们做不到的事情：选定一组让所有惯性观者都认同为粒子的、自然的基模式。在 9.2 节的末尾，我们简要讨论过受到瞬态外力作用的振子，以及如何定义一个 S 矩阵，把早期时刻的粒子数本征态与晚期时刻的粒子数本征态联系起来。同一套想法可以直接搬到量子场论中。如果我们考虑这样一种情形：时空在渐近过去和渐近未来都是静态的，而中间存在某种扰动，那么我们就可以定义 in 态与 out 态——它们分别是早期和晚期的能量本征态——以及一组 Bogoliubov 系数，描述 in 真空（比如说）用 out 态表出时会呈现为怎样的多粒子位形。这种现象被称为引力场导致的粒子产生（particle production by gravitational fields）\footnote{有意思的是，关于弯曲时空中粒子产生的最早讨论正是 Schrödinger 本人给出的；见 E. Schrödinger (1939), \emph{Physica} (Utrecht) 6, 899.}；相关的物理例子包括早期宇宙和黑洞。

\section{安鲁效应}

必须承认，尽管我们在理解弯曲时空量子场论的基础上下了一番功夫，实际上却不会在任何弯曲背景中做任何详细的计算。作为替代，我们将研究一种现象，它依赖我们已引入的那些观念，却甚至在平直时空中就会显现：这就是安鲁效应（Unruh effect），它断言处于传统 Minkowski 真空态中的加速观者会观测到粒子的热谱。历史上，安鲁效应正是在试图理解霍金效应（黑洞事件视界存在时的热辐射）背后的物理时被发现的。我们的策略将是仔细推导安鲁效应，然后在下一节中在合理假设下论证它蕴含霍金效应；霍金效应难以直接推导，仅仅因为求解弯曲时空中的波动方程比在平直时空中更难。

安鲁效应的基本想法很简单：它是下述观念的一个体现——对正频与负频模式持有不同理解的观者，对同一个态的粒子内容也会有不同的看法。对于 Minkowski 空间中匀加速运动的观者，其轨迹将沿一条类时 Killing 矢量的轨道行进，但沿的不是通常时间平移对称性的轨道。因此我们可以用适合加速观者的模式来展开场，并在通常的 Minkowski 真空中计算粒子数算符，届时会发现一个粒子热谱。对这个结果可以配以不同套的解释性说法；应当汲取的基本教训是：我们心目中毫无动静的真空，实际上带有热态的特征。

为了剔除一切可能的复杂因素、直抵底层现象，我们考虑一个在不至于完全平庸的前提下尽可能简单的量子场论：两维时空（$n=2$）中的无质量（$m=0$）标量场。在两维中，共形耦合与最小耦合相重合，所以我们不包含与标量曲率的任何直接相互作用。（反正我们身处平直时空，这种耦合本来也不会有任何效果。）相应的波动方程为

\begin{equation*}
\Box\phi = 0.
\tag{9.124}
\end{equation*}

在深入这个场论的量子化之前，我们先来想一想一位匀加速观者眼中的两维 Minkowski 空间。我们知道，度规在惯性坐标下可以写成

\begin{equation*}
ds^2 = -dt^2 + dx^2.
\tag{9.125}
\end{equation*}

考虑一位沿 $x$ 方向以大小为 $\alpha$ 的均匀加速度运动的观者。我们宣称，由此得到的轨迹 $x^{\mu}(\tau)$ 将由下式给出：

\begin{align*}
t(\tau) &= \frac{1}{\alpha}\sinh(\alpha\tau) \\
x(\tau) &= \frac{1}{\alpha}\cosh(\alpha\tau).
\tag{9.126}
\end{align*}

我们来验证这条路径确实对应于恒定的加速度。加速度这个二矢量在整体惯性坐标系中由下式给出：

\begin{equation*}
a^{\mu} = \frac{D^2 x^{\mu}}{d\tau^2} = \frac{d^2 x^{\mu}}{d\tau^2},
\tag{9.127}
\end{equation*}

其中沿路径的协变导数等于普通导数，因为 Christoffel 符号在这些坐标下为零。于是 $a^{\mu}$ 的分量为

\begin{align*}
a^t &= \alpha\sinh(\alpha\tau) \\
a^x &= \alpha\cosh(\alpha\tau),
\tag{9.128}
\end{align*}

而大小为

\begin{equation*}
\sqrt{a_{\mu}a^{\mu}} = \sqrt{-\alpha^{2}\sinh^{2}(\alpha\tau) + \alpha^{2}\cosh^{2}(\alpha\tau)} = \alpha.
\tag{9.129}
\end{equation*}

因此这条路径确实对应于大小为 $\alpha$ 的恒定加速度，正如所愿。我们这位加速观者的轨迹满足关系

\begin{equation*}
x^2(\tau) = t^2(\tau) + 1/\alpha^2,
\tag{9.130}
\end{equation*}

从而描述一个双曲面，它在过去渐近于类光路径 $x=-t$，在未来渐近于 $x=t$。加速观者从过去类光无限远行进到未来类光无限远，而不是像测地观者那样到达类时无限远。

我们可以在两维 Minkowski 空间上选取适配匀加速运动的新坐标 $(\eta,\xi)$。令

\begin{equation*}
t = \frac{1}{a}e^{a\xi}\sinh(a\eta), \qquad x = \frac{1}{a}e^{a\xi}\cosh(a\eta) \qquad (x > |t|).
\tag{9.131}
\end{equation*}

\begin{figure}[H]
\centering
\includegraphics[width=0.72\textwidth]{fig_9.1.pdf}
\caption*{图 9.1\quad Rindler 坐标下的 Minkowski 时空。区域 I 是沿 $+x$ 方向做恒定加速的观者所能到达的区域。坐标 $(\eta,\xi)$ 可用于区域 I，也可单独用于区域 IV，但此时它们指向相反的方向。矢量场 $\partial_{\eta}$ 对应于洛伦兹推动（boost）对称性的生成元。视界 $H^{\pm}$ 是该矢量场的 Killing 视界，同时也是 Rindler 观者所见过去与未来的边界。}
\end{figure}

新坐标的取值范围为

\begin{equation*}
-\infty < \eta,\, \xi < +\infty,
\tag{9.132}
\end{equation*}

并且覆盖楔形区域 $x>|t|$，即图 9.1 中标记为区域 I 的部分。在这些坐标中，恒定加速路径 (9.126) 表为

\begin{align*}
\eta(\tau) &= \frac{\alpha}{a}\tau \\
\xi(\tau) &= \frac{1}{a}\ln\left(\frac{a}{\alpha}\right),
\tag{9.133}
\end{align*}

于是固有时正比于 $\eta$，而空间坐标 $\xi$ 为常数。特别地，$\alpha=a$ 的观者沿如下路径运动：

\begin{equation*}
\eta = \tau, \qquad \xi = 0.
\tag{9.134}
\end{equation*}

在这些坐标中度规取如下形式：

\begin{equation*}
ds^2 = e^{2a\xi}(-d\eta^2 + d\xi^2).
\tag{9.135}
\end{equation*}

带着这个度规，区域 I 被称为\textbf{Rindler 空间}（Rindler space），尽管它显然只是 Minkowski 空间的一部分。\textbf{Rindler 观者}（Rindler observer）就是沿恒定加速路径运动的观者，如 (9.133) 那样。Rindler 空间的因果结构酷似图 5.12 中 Schwarzschild 解最大延拓的 $r>2GM$ 区域。特别地，图 9.1 中标记为 $H^{+}$ 的类光直线 $x=t$，对于区域 I 中任何 $\eta=$ 常数的类空超曲面而言都是未来 Cauchy 视界；类似地，$H^{-}$ 是过去 Cauchy 视界。这些视界让人想起 Kruskal 图中的事件视界：Schwarzschild 中的静态观者（$r=$ 常数）正对应于 Rindler 空间中的恒定加速路径。

(9.135) 中的度规分量与 $\eta$ 无关，所以我们立刻知道 $\partial_{\eta}$ 是一个 Killing 矢量。但当然，这本来就是 Minkowski 时空，所以我们自以为了解所有 Killing 矢量是什么。确实，如果把 $\partial_{\eta}$ 用 $(t,x)$ 坐标表出，我们发现

\begin{align*}
\partial_{\eta} &= \frac{\partial t}{\partial\eta}\partial_t + \frac{\partial x}{\partial\eta}\partial_x \\
&= e^{a\xi}\left[\cosh(a\eta)\partial_t + \sinh(a\eta)\partial_x\right] \\
&= a(x\partial_t + t\partial_x).
\tag{9.136}
\end{align*}

这不多不少，正是与沿 $x$ 方向的推动（boost）相联系的那个 Killing 场。从这个表达式可以清楚地看到，这个 Killing 场自然地延拓到整个时空；在区域 II 和 III 中它是类空的，而在区域 IV 中它是类时的但指向过去。我们上面辨认出的那些视界，实际上就是 $\partial_{\eta}$ 的 Killing 视界。红移因子（在式 (6.12) 中定义为 Killing 矢量范数的大小）为

\begin{equation*}
V = e^{a\xi}.
\tag{9.137}
\end{equation*}

这个 Killing 视界的表面引力 $\kappa=\sqrt{\nabla_{\mu}V\nabla^{\mu}V}$ 于是为

\begin{equation*}
\kappa = a.
\tag{9.138}
\end{equation*}

这里并没有什么真实的引力，因为我们身处平直空间；但这个表面引力刻画了 Rindler 观者的加速度。

我们还可以把 (9.131) 中的符号翻转，在区域 IV 中定义坐标 $(\eta,\xi)$：

\begin{equation*}
t = -\frac{1}{a}e^{a\xi}\sinh(a\eta), \qquad x = -\frac{1}{a}e^{a\xi}\cosh(a\eta) \qquad (x < |t|).
\tag{9.139}
\end{equation*}

这个符号保证了在区域 IV 中 $\partial_{\eta}$ 与 $\partial_t$ 指向相反的方向。严格说来，我们不能在区域 I 和区域 IV 中同时使用 $(\eta,\xi)$，因为这两个坐标在各自区域中的取值范围相同；不过只要我们明确指明所谈论的是哪个区域，就不会出问题。之所以宁愿把同一套坐标标签用两次，而不是干脆引入新坐标，是因为度规 (9.135) 对区域 I 和区域 IV 都适用。

在曲面 $t=0$ 上，$\partial_{\eta}$ 是一个超曲面正交的类时 Killing 矢量，唯一的例外是它为零的那一点 $x=0$。因此这个矢量可用来定义一组正频与负频模式，并在此基础上为标量场的希尔伯特空间搭建 Fock 基。Rindler 坐标下无质量 Klein--Gordon 方程的形式为

\begin{equation*}
\Box\phi = e^{-2a\xi}\left(-\partial_{\eta}^{2} + \partial_{\xi}^{2}\right)\phi = 0.
\tag{9.140}
\end{equation*}

归一化的平面波 $g_k = (4\pi\omega)^{-1/2}e^{-i\omega\eta+ik\xi}$（其中 $\omega=|k|$）求解这个方程，而且显然具有正频，即 $\partial_{\eta}g_k=-i\omega g_k$。但我们需要的模式要相对于一个指向未来的 Killing 矢量具有正频，而在区域 IV 中扮演这一角色的是 $\partial_{(-\eta)}=-\partial_{\eta}$ 而非 $\partial_{\eta}$。为了对付这个麻烦，我们引入两组模式，一组支撑在区域 I 中，另一组支撑在区域 IV 中：

\begin{equation*}
\begin{aligned}
g_{k}^{(1)} &= \begin{cases} \dfrac{1}{\sqrt{4\pi\omega}}\, e^{-i\omega\eta+ik\xi} & \mathrm{I} \\[10pt] 0 & \mathrm{IV} \end{cases} \\[14pt]
g_{k}^{(2)} &= \begin{cases} 0 & \mathrm{I} \\[10pt] \dfrac{1}{\sqrt{4\pi\omega}}\, e^{+i\omega\eta+ik\xi} & \mathrm{IV} \end{cases}
\end{aligned}
\tag{9.141}
\end{equation*}

两种情形下都取 $\omega=|k|$；在两维中，空间波矢就是单独一个数 $k$。每一组模式都相对于相应的指向未来的类时 Killing 矢量具有正频：

\begin{gather*}
\partial_{\eta}\, g_{k}^{(1)} = -i\omega g_{k}^{(1)} \\
\partial_{(-\eta)}\, g_{k}^{(2)} = -i\omega g_{k}^{(2)}, \qquad \omega > 0.
\tag{9.142}
\end{gather*}

这两组模式连同它们的共轭，构成了对整个时空中波动方程任何解都完备的一组基模式。（单独一点 $x=t=0$ 是测度为零的集合，用不着为此担心。）在 Rindler 图的区域 II 和 III 中，这两组模式都不为零；把它们写成坐标 $\eta$ 和 $\xi$ 的形式时这一点被掩盖了，但这些函数可以解析延拓到未来和过去的区域。把相应的湮灭算符记作 $\hat{b}_{k}^{(1,2)}$，我们可以写出

\begin{equation*}
\phi = \int dk\left(\hat{b}_{k}^{(1)}g_{k}^{(1)} + \hat{b}_{k}^{(1)\dagger}g_{k}^{(1)*} + \hat{b}_{k}^{(2)}g_{k}^{(2)} + \hat{b}_{k}^{(2)\dagger}g_{k}^{(2)*}\right).
\tag{9.143}
\end{equation*}

这一展开是 (9.66) 式——用原始 Minkowski 模式所作的展开——的另一种选择；在两维中该式取如下形式：

\begin{equation*}
\phi = \int dk\left(\hat{a}_{k}f_{k} + \hat{a}_{k}^{\dagger}f_{k}^{*}\right).
\tag{9.144}
\end{equation*}

容易直接检验，模式 (9.141) 相对于内积 (9.94) 是恰当归一化的。在度规 (9.135) 中，曲面 $\eta=0$ 的指向未来的单位法矢量归一化为

\begin{equation*}
-1 = g_{\mu\nu}n^{\mu}n^{\nu} = -e^{2a\xi}(n^{0})^{2},
\tag{9.145}
\end{equation*}

即

\begin{equation*}
n^{0} = e^{-a\xi}.
\tag{9.146}
\end{equation*}

与此同时，空间度规的行列式满足

\begin{equation*}
\sqrt{\gamma} = e^{a\xi}.
\tag{9.147}
\end{equation*}

于是我们有 $n^{0}\sqrt{\gamma}=1$，而 Rindler 模式内积的计算与普通 Minkowski 模式的计算如出一辙。最终得到

\begin{align*}
(g_{k_1}^{(1)}, g_{k_2}^{(1)}) &= \delta(k_1 - k_2) \\
(g_{k_1}^{(2)}, g_{k_2}^{(2)}) &= \delta(k_1 - k_2) \\
(g_{k_1}^{(1)}, g_{k_2}^{(2)}) &= 0,
\tag{9.148}
\end{align*}

共轭模式也有类似结果。

这样就有两组模式——Minkowski 模式与 Rindler 模式——都可以用来展开平直两维时空中 Klein--Gordon 方程的解。虽然理论的希尔伯特空间在两种表示下是同一个，但它作为 Fock 空间的诠释却会不同；特别地，真空态会不相同。Minkowski 真空 $|0_{\mathrm{M}}\rangle$ 满足

\begin{equation*}
\hat{a}_{k}|0_{\mathrm{M}}\rangle = 0,
\tag{9.149}
\end{equation*}

在 Rindler 表示中它将被描述为一个多粒子态；同样，Rindler 真空 $|0_{\mathrm{R}}\rangle$ 满足

\begin{equation*}
\hat{b}_{k}^{(1)}|0_{\mathrm{R}}\rangle = \hat{b}_{k}^{(2)}|0_{\mathrm{R}}\rangle = 0,
\tag{9.150}
\end{equation*}

在 Minkowski 表示中它将被描述为一个多粒子态。在实际层面上，这种差异源于：单个 Rindler 模式永远无法写成正频 Minkowski 模式之和；在 $t=0$ 处 Rindler 模式只有支撑在半直线上的部分，而这样的函数不可能展开成纯粹正频的平面波。因此，用来定义 $|0_{\mathrm{R}}\rangle$ 的那些 Rindler 湮灭算符必然是 Minkowski 产生算符与湮灭算符的叠加，所以两个真空不可能重合。

Rindler 观者相对于推动 Killing 矢量 $\partial_{\eta}$ 的轨道是静止的。因此区域 I 中的这种观者将用 Rindler 模式 $g_{k}^{(1)}$ 来描述粒子：特别地，他会观测到 Rindler 真空中的态不含粒子，态 $\hat{b}_{k}^{(1)\dagger}|0_{\mathrm{R}}\rangle$ 含有一个频率为 $\omega=|k|$ 的粒子，如此等等。反过来，一位穿过 Minkowski 真空态的 Rindler 观者将探测到一个粒子背景，尽管惯性观者会把这个态描述成完全空无一物。Rindler 观者会探测到什么样的粒子呢？我们知道怎样回答这个问题：计算联系 Minkowski 模式与 Rindler 模式的 Bogoliubov 系数，再用它们确定 Rindler 粒子数算符在 Minkowski 真空中的期望值。这直截了当却十分繁琐，所以我们将走一条由 Unruh 提出的捷径。我们将找到一组模式，它与 Minkowski 模式共享同一个真空态（尽管对激发态的描述可能不同），但它与 Rindler 模式的交叠更为直接。做法是：从 Rindler 模式出发，把它们解析延拓到整个时空，再把这一延拓用原始的 Rindler 模式表出。

为了看清这是如何运作的，注意由 (9.131) 和 (9.139) 可知，在区域 I 和 IV 中，Minkowski 坐标 $(t,x)$ 与 Rindler 坐标 $(\eta,\xi)$ 之间有如下关系：

\begin{equation*}
\begin{aligned}
e^{-a(\eta-\xi)} &= \begin{cases} a(-t+x) & \mathrm{I} \\ a(t-x) & \mathrm{IV} \end{cases} \\[10pt]
e^{a(\eta+\xi)} &= \begin{cases} a(t+x) & \mathrm{I} \\ a(-t-x) & \mathrm{IV} \end{cases}
\end{aligned}
\tag{9.151}
\end{equation*}

因此，对于 $k>0$（从而 $\omega=k$）的模式 $g_{k}^{(1)}$，其在区域 I 中用 Minkowski 坐标表出的时空依赖为

\begin{align*}
\sqrt{4\pi\omega}\, g_{k}^{(1)} &= e^{-i\omega\eta+ik\xi} \\
&= e^{-i\omega(\eta-\xi)} \\
&= a^{i\omega/a}(-t+x)^{i\omega/a}.
\tag{9.152}
\end{align*}

这个函数在整个时空中的解析延拓很直截了当：对任何 $(t,x)$ 值都用这最后一个表达式即可。但我们希望处处都用原始的 Rindler 模式来表示结果；由于 $g_{k}^{(1)}$ 模式在区域 IV 中为零，我们需要让 $g_{k}^{(2)}$ 模式登场。在区域 IV 中用 Minkowski 坐标表出这些模式，对 $k>0$ 得到

\begin{align*}
\sqrt{4\pi\omega}\, g_{k}^{(2)} &= e^{+i\omega\eta+ik\xi} \\
&= e^{+i\omega(\eta+\xi)} \\
&= a^{-i\omega/a}(-t-x)^{-i\omega/a}.
\tag{9.153}
\end{align*}

这与我们想要的 (9.152) 的行为并不匹配。但若取复共轭并反转波数，就得到

\begin{align*}
\sqrt{4\pi\omega}\, g_{-k}^{(2)*} &= e^{-i\omega\eta+ik\xi} \\
&= e^{-i\omega(\eta-\xi)} \\
&= a^{i\omega/a}(t-x)^{i\omega/a} \\
&= a^{i\omega/a}[e^{-i\pi}(-t+x)]^{i\omega/a} \\
&= a^{i\omega/a}e^{\pi\omega/a}(-t+x)^{i\omega/a}.
\tag{9.154}
\end{align*}

组合

\begin{equation*}
\sqrt{4\pi\omega}\left(g_{k}^{(1)} + e^{-\pi\omega/a}g_{-k}^{(2)*}\right) = a^{i\omega/a}(-t+x)^{i\omega/a}
\tag{9.155}
\end{equation*}

因此沿整个曲面 $t=0$ 都有良好定义。我们明确考察了 $k>0$ 的情形，但对 $k<0$ 会得到完全相同的结果。

这个模式的恰当归一化版本由下式给出：

\begin{equation*}
h_{k}^{(1)} = \frac{1}{\sqrt{2\sinh\left(\frac{\pi\omega}{a}\right)}}\left(e^{\pi\omega/2a}g_{k}^{(1)} + e^{-\pi\omega/2a}g_{-k}^{(2)*}\right).
\tag{9.156}
\end{equation*}

这是 $g_{k}^{(1)}$ 模式的一个恰当的解析延拓；为了得到完备集，还需纳入 $g_{k}^{(2)}$ 模式的延拓，由类似的论证可知它为

\begin{equation*}
h_{k}^{(2)} = \frac{1}{\sqrt{2\sinh\left(\frac{\pi\omega}{a}\right)}}\left(e^{\pi\omega/2a}g_{k}^{(2)} + e^{-\pi\omega/2a}g_{-k}^{(1)*}\right).
\tag{9.157}
\end{equation*}

为了验证归一化——比如对 $h_{k}^{(1)}$——我们利用 (9.148)：

\begin{align*}
\left(h_{k_1}^{(1)}, h_{k_2}^{(1)}\right) &= \frac{1}{2\sqrt{\sinh\left(\frac{\pi\omega_1}{a}\right)\sinh\left(\frac{\pi\omega_2}{a}\right)}}\Bigg[e^{\pi(\omega_1+\omega_2)/2a}\left(g_{k_1}^{(1)}, g_{k_2}^{(1)}\right) \\
&\quad + e^{-\pi(\omega_1+\omega_2)/2a}\left(g_{-k_1}^{(2)*}, g_{-k_2}^{(2)*}\right)\Bigg] \\
&= \frac{1}{2\sqrt{\sinh\left(\frac{\pi\omega_1}{a}\right)\sinh\left(\frac{\pi\omega_2}{a}\right)}}\Bigg[e^{\pi(\omega_1+\omega_2)/2a}\,\delta(k_1 - k_2) \\
&\quad + e^{-\pi(\omega_1+\omega_2)/2a}\,\delta(-k_1 + k_2)\Bigg] \\
&= \frac{e^{\pi\omega_1/a} - e^{-\pi\omega_1/a}}{2\sinh\left(\frac{\pi\omega_1}{a}\right)}\,\delta(k_1 - k_2) \\
&= \delta(k_1 - k_2),
\tag{9.158}
\end{align*}

正如我们所愿。

现在可以把场用这些模式展开：

\begin{equation*}
\phi = \int dk\left(\hat{c}_{k}^{(1)}h_{k}^{(1)} + \hat{c}_{k}^{(1)\dagger}h_{k}^{(1)*} + \hat{c}_{k}^{(2)}h_{k}^{(2)} + \hat{c}_{k}^{(2)\dagger}h_{k}^{(2)*}\right).
\tag{9.159}
\end{equation*}

根据 9.4 节中关于 Bogoliubov 变换的讨论我们知道，用 $g_{k}^{(1,2)}$ 模式表出 $h_{k}^{(1,2)}$ 模式的表达式 (9.156) 与 (9.157)，蕴含着用算符 $\hat{c}_{k}^{(1,2)}$ 表出 Rindler 算符 $\hat{b}_{k}^{(1,2)}$ 的相应表达式：

\begin{align*}
\hat{b}_{k}^{(1)} &= \frac{1}{\sqrt{2\sinh\left(\frac{\pi\omega}{a}\right)}}\left(e^{\pi\omega/2a}\hat{c}_{k}^{(1)} + e^{-\pi\omega/2a}\hat{c}_{-k}^{(2)\dagger}\right) \\
\hat{b}_{k}^{(2)} &= \frac{1}{\sqrt{2\sinh\left(\frac{\pi\omega}{a}\right)}}\left(e^{\pi\omega/2a}\hat{c}_{k}^{(2)} + e^{-\pi\omega/2a}\hat{c}_{-k}^{(1)\dagger}\right).
\tag{9.160}
\end{align*}

于是我们可以把区域 I 中的 Rindler 粒子数算符

\begin{equation*}
\hat{n}_{\mathrm{R}}^{(1)}(k) = \hat{b}_{k}^{(1)\dagger}\hat{b}_{k}^{(1)},
\tag{9.161}
\end{equation*}

用新的算符 $\hat{c}_{k}^{(1,2)}$ 表示。

原先 $k>0$ 的正频 Minkowski 平面波模式 $f_k\propto e^{-i\omega(t-x)}$，只要 $\mathrm{Im}(t-x)\leq 0$，对复的 $(t,x)$ 都是解析且有界的。（这类模式称为“右行”（right-moving），因为它们描述向右传播的波。）只要我们把虚幂函数的支割线取在上半复 $(t-x)$ 平面内，我们的新模式 $h_{k}^{(1)}$ 也具有同样的性质——考察 (9.152) 和 (9.154) 即可看出；这与我们在 (9.154) 中取 $-1=e^{-i\pi}$ 的做法是一致的。类似的考虑也适用于 $h_{k}^{(2)}$ 模式，它们在下半复 $(t+x)$ 平面内解析且有界，正如 $k<0$ 的正频 Minkowski 平面波模式（左行）一样。因此，与原先的 Rindler 模式 $g_{k}^{(1,2)}$ 不同，我们知道 $h_{k}^{(1,2)}$ 模式能够纯粹用正频 Minkowski 模式 $f_k$ 表出。所以它们与 Minkowski 模式共享同一个真空态 $|0_{\mathrm{M}}\rangle$，于是

\begin{equation*}
\hat{c}_{k}^{(1)}|0_{\mathrm{M}}\rangle = \hat{c}_{k}^{(2)}|0_{\mathrm{M}}\rangle = 0.
\tag{9.162}
\end{equation*}

激发态将不会重合，但这不会困扰我们，因为我们关心的正是当态恰好处于 Minkowski 真空时 Rindler 观者会看到什么。例如，区域 I 中的观者将观测由算符 $\hat{b}_{k}^{(1)}$ 定义的粒子；频率为 $\omega$ 的这类粒子的期望数目由下式给出：

\begin{align*}
\langle 0_{\mathrm{M}}|\hat{n}_{\mathrm{R}}^{(1)}(k)|0_{\mathrm{M}}\rangle &= \langle 0_{\mathrm{M}}|\hat{b}_{k}^{(1)\dagger}\hat{b}_{k}^{(1)}|0_{\mathrm{M}}\rangle \\
&= \frac{1}{2\sinh\left(\frac{\pi\omega}{a}\right)}\langle 0_{\mathrm{M}}|e^{-\pi\omega/a}\hat{c}_{-k}^{(1)}\hat{c}_{-k}^{(1)\dagger}|0_{\mathrm{M}}\rangle \\
&= \frac{e^{-\pi\omega/a}}{2\sinh\left(\frac{\pi\omega}{a}\right)}\,\delta(0) \\
&= \frac{1}{e^{2\pi\omega/a} - 1}\,\delta(0),
\tag{9.163}
\end{align*}

其中我们用到了 $\hat{c}_{k}^{(1)\dagger}|0_{\mathrm{M}}\rangle$ 是归一化的单粒子态这一事实：

\begin{equation*}
\langle 0_{\mathrm{M}}|\hat{c}_{k}^{(1)}\hat{c}_{k}^{(1)\dagger}|0_{\mathrm{M}}\rangle = \delta(0).
\tag{9.164}
\end{equation*}

(9.163) 中的 $\delta$ 函数只是我们使用（不可平方可积的）平面波基模式带来的假象；如果我们构造的是归一化的波包，就会得到一个有限的结果，而谱完全相同。

结果 (9.163) 是一个温度为

\begin{equation*}
T = \frac{a}{2\pi}.
\tag{9.165}
\end{equation*}

的 Planck 谱。这样，\emph{以均匀加速度穿过 Minkowski 真空运动的观者会观测到粒子的热谱}。这就是\textbf{安鲁效应}。当然，热辐射不只是谱 (9.163) 而已；要真正称得上热，我们还应当检查在观测到的粒子之间没有隐匿的关联。这一点已经得到验证：Rindler 观者探测到的辐射是真正热的。在最基本的层面上，安鲁效应表明两批不同的观者（惯性的与 Rindler 的）会用极为不同的语言描述同一个态；在稍稍深入的层面上，它揭示了量子场论中真空本质上的热性质。

温度 $T=a/2\pi$ 是沿路径 $\xi=0$ 运动的观者会测得的温度，这条路径感受到的加速度为 $\alpha=a$。利用 (9.133) 可知，任何 $\xi=$ 常数的其他路径感受到的加速度为

\begin{equation*}
\alpha = ae^{-a\xi}
\tag{9.166}
\end{equation*}

因而应当测得温度为 $\alpha/2\pi$ 的热辐射。这与我们在第 6 章中关于沿某个 Killing 矢量 $K^{\mu}$ 的轨道运动的静态观者所见证的红移的讨论相一致；我们当时发现，在点 $x_1$ 处以频率 $\omega_1$ 发出的辐射，将在点 $x_2$ 处被观测到具有频率

\begin{equation*}
\omega_2 = \frac{V_1}{V_2}\omega_1,
\tag{9.167}
\end{equation*}

其中红移因子 $V$ 是 Killing 矢量的范数。在 (9.137) 中我们发现与 $\partial_{\eta}$ 相联系的红移因子为 $V=e^{a\xi}$，于是

\begin{equation*}
\omega_2 = e^{a(\xi_1 - \xi_2)}\omega_1.
\tag{9.168}
\end{equation*}

因此，如果 $\xi_1=0$ 处的观者探测到温度 $T=a/2\pi$，那么 $\xi_2=\xi$ 处的观者将看到它红移为温度 $T=ae^{-a\xi}/2\pi$，正如式 (9.166) 那样。特别地，当 $\xi\to+\infty$ 时温度一路红移到零。这是有道理的：无穷远处的 Rindler 观者近乎惯性观者，会对真空和粒子给出与普通 Minkowski 观者相同的定义。

安鲁效应告诉我们，加速观者会在 Minkowski 真空态中探测到粒子。当然，惯性观者会把同一个态描述成完全空无一物；实际上，能动张量的期望值会是 $\langle T_{\mu\nu}\rangle=0$。可是如果没有能量-动量，Rindler 观者们怎么能探测到粒子呢？这是一个微妙的问题，但绝不是矛盾。Rindler 观者要探测背景粒子，就必须携带一台探测器——某种与被探测粒子相耦合的装置。但如果探测器被维持在恒定加速度下，能量就不守恒；我们需要持续对探测器做功才能让它一直加速。从 Minkowski 观者的角度看，Rindler 探测器\emph{发射}粒子，也同样吸收粒子；一旦引入耦合，发射的可能性就不可避免。当探测器记录到一个粒子时，惯性观者会说它发射了一个粒子，并因此感受到一个辐射反作用力。归根结底，激发 Rindler 探测器所需的能量并非来自背景能动张量，而是来自我们为维持探测器加速而注入的能量。

\section{霍金效应与黑洞蒸发}

尽管发生在平直时空，安鲁效应却给了我们弯曲时空量子场论最重要的一课：“真空”和“粒子”是依赖于观者的概念，而非基本概念。实际上，鉴于我们对安鲁效应的理解，几乎可以立刻看出霍金效应是如何产生的。这不应太令人惊讶，因为我们已经指出过 Rindler 空间的因果结构与描述永恒黑洞的最大延拓 Schwarzschild 时空的因果结构之间的相似性。因此，我们将能够论证霍金辐射的存在，而无需在弯曲时空中做任何显式计算；当然，还有许多你可能想更深入考察的特征，那就需要充分动用弯曲度规的力量。除 Birrell 和 Davies (1982) 与 Wald (1994) 之外，还有一些很好的综述文章，你可以在其中找到对这里讨论的问题更全面的讨论。\footnote{T.A. Jacobson, “Introductory Lectures on Black Hole Thermodynamics,” Lectures at University of Utrecht (1996), http://www.fys.ruu.nl/\textasciitilde{}wwwthe/lectures/itfuu-0196.ps; R.M. Wald, “The thermodynamics of black holes,” \emph{Living Rev.\ Rel.} \textbf{4}, 6 (2001), http://arxiv.org/gr-qc/9912119; J. Traschen, “An introduction to black hole evaporation” (2000), http://arxiv.org/gr-qc/0010055.} 我们对霍金辐射的推导沿循 Jacobson 的做法。

考虑一个位于 Schwarzschild 黑洞之外、半径 $r_{1} > 2GM$ 处的静态观者。这样的观者沿着类时 Killing 矢量 $K = \partial_{t}$ 的轨道运动。在第 6 章中我们已经证明，对 Schwarzschild 时空中的静态观者，红移因子 $V = \sqrt{-K_{\mu}K^{\mu}}$ 由下式给出

\begin{equation*}
V = \sqrt{1 - \frac{2GM}{r}},
\tag{9.169}
\end{equation*}

相应的加速度大小则由下式给出：

\begin{equation*}
a = \frac{GM}{r\sqrt{r - 2GM}}.
\tag{9.170}
\end{equation*}

对于非常靠近事件视界的观者，$r_{1} - 2GM \ll 2GM$，这一加速度与 Schwarzschild 半径所标定的尺度相比会变得非常大，

\begin{equation*}
a_{1} \gg \frac{1}{2GM}.
\tag{9.171}
\end{equation*}

Schwarzschild 半径又标定了视界附近时空的曲率半径。因此，在由 $a_{1}^{-1} \ll 2GM$ 所设定的长度和时间尺度上观察，时空看上去基本上是平直的。让我们作一个至关重要的假设：在黑洞附近\emph{自由下落}的观者看来，某个标量场 $\phi$ 的量子态就像 Minkowski 真空（不含任何粒子）。这个假设是合理的，因为事件视界并不是一种局域的屏障；自由下落的观者在穿过视界时看不到任何异常发生。于是，这位静态观者看起来就跟平直时空中的常加速度观者一模一样，将探测到温度为 $T_{1} = a_{1}/2\pi$ 的安鲁辐射。

现在考虑位于无限远处——或者至少位于与 $2GM$ 相比大得多的距离 $r_{2}$ 处——的静态观者。在这种情形下，无论如何都谈不上能在 $a_{2}^{-1} \gg 2GM$ 的时间尺度上忽略时空曲率，因此没有理由预期他们会看到温度为 $a_{2}/2\pi$ 的辐射，这里 $a_{2}$ 在 $r_{2}$ 处取值。但是在视界附近观测到的辐射会带着恰当的红移传播到无限远。我们可以套用上一节末尾用过的论证，来确定这样的观者应当会看到什么；他们应当探测到红移至如下温度的热辐射：

\begin{equation*}
T_{2} = \frac{V_{1}}{V_{2}} T_{1} = \frac{V_{1}}{V_{2}} \frac{a}{2\pi}.
\tag{9.172}
\end{equation*}

在无限远处 $V_{2} \to 1$，于是观测到的温度为

\begin{equation*}
T = \lim_{r_{1} \to 2GM} \frac{V_{1} a_{1}}{2\pi} = \frac{\kappa}{2\pi},
\tag{9.173}
\end{equation*}

其中 $\kappa = \lim(Va)$ 是表面引力；对 Schwarzschild 黑洞，$\kappa = 1/4GM$。与平直时空中的加速观者不同，Schwarzschild 时空中的静态 Killing 矢量在无限远处具有有限的模长，因此视界附近的辐射只红移到一个有限值，而不会一直红移到零。远离黑洞的观者因此会看到黑洞发出的热辐射流，其温度正比于黑洞的表面引力。这就是著名的\textbf{霍金效应}（Hawking effect），这种辐射本身则被称为霍金辐射。

这种对霍金效应的推导虽然取巧，却没有任何不诚实之处。特别是，与加速度的联系清楚地说明了为什么温度正比于黑洞的表面引力（这一结论对更一般的黑洞同样成立，而不只是 Schwarzschild 黑洞）。不过，我们需要对我们所作的那个假设有清醒的认识：视界附近的真空态在自由下落的观者看来是非奇异的。用术语来说，就是取重整化能动张量在视界处为有限，或者等价地，两点函数满足 Hadamard 条件 (9.123)。

通过考察最大延拓 Schwarzschild 几何中各种可能的真空态，这个假设的含义会变得更加清楚。这些态对于由引力坍缩形成的现实黑洞来说不一定具有物理上的相关性，但在这种理想化情形中出现的各种可能性，能给现实世界带来有益的教益。我们将只描述这些态，而不定量地刻画它们，也不推导它们的任何性质；更多细节可参见上文所引的文献。

在寻找真空态时，我们也许会先去寻找一个在整个时空中都正则的态［在 Hadamard 意义下，式 (9.123)］。对最大延拓的 Schwarzschild 时空，Hartle 和 Hawking 找到了这样一个态，我们称之为\textbf{Hartle--Hawking 真空}（Hartle--Hawking vacuum）；的确，这是唯一的既处处正则、又在代表无限远处时间平移的 Schwarzschild Killing 矢量 $\partial_{t}$ 作用下不变的真空态。特别是，回忆图 5.16 所示的 Schwarzschild 共形图，Hartle--Hawking 真空在 $r = 2GM$ 处的过去和未来事件视界 $H^{\pm}$ 上正则，在过去和未来类光无限远 $\mathscr{I}^{\pm}$ 上也正则。根据上文对静态观者的讨论，我们应当会预期 Hartle--Hawking 真空中有热辐射从黑洞发出，而事实证明确实如此。然而，对这个态的仔细考察揭示出，还有一份等量的热辐射流从过去类光无限远（$\mathscr{I}^{-}$）流向黑洞；换句话说，它描述的是一个与其环境处于热平衡的黑洞。我们并不会用它来模拟我们宇宙中现实的黑洞。另一种真空更类似于经由引力坍缩形成的黑洞的真空态，这就是\textbf{Unruh 真空}（Unruh vacuum），它在 $H^{+}$ 上非奇异（因此预言了出射的霍金辐射），但没有来自 $\mathscr{I}^{-}$ 的入射辐射。Unruh 真空在 Schwarzschild 的过去视界 $H^{-}$ 上是奇异的；如果我们只是把它当作现实黑洞的模型，这并不会困扰我们，因为如图 5.17 那样经历引力坍缩的时空不会有白洞，也没有任何过去视界。最后，我们也许会去寻找这样一种真空态：既没有粒子进入黑洞，也没有粒子逃往无限远；换句话说，$\mathscr{I}^{\pm}$ 上通量为零。确实存在这样一个态，称为\textbf{Boulware 真空}（Boulware vacuum）。这样一种态的存在，似乎与我们从安鲁效应出发对霍金效应所作的论证相冲突，但仔细的分析表明，Boulware 真空在 $H^{-}$ 和 $H^{+}$ 上都是奇异的。于是，认为真空在视界附近的自由下落观者看来是正则的这一假设，在这个态中被违反了。

因此，对永恒 Schwarzschild 度规中真空态的仔细考察，与我们基于安鲁效应的推理是一致的；在 $H^{+}$ 上正则的态预言了预期形式的霍金辐射。注意，事件视界的存在对这一论证至关重要；如果没有这样的视界，态必须在视界上正则的要求就没有任何约束力。以中子星为例：其半径可能接近 Schwarzschild 半径，但其时空没有任何视界。中子星不发出任何霍金辐射。理解这一点的一个途径是认识到，静态中子星度规具有一个处处类时的 Killing 矢量，可以用它来定义贯穿整个时空延伸、并且在无限远处与 Minkowski 模式相匹配的正频模式。这样得到的真空态实际上会类似于 Boulware 真空，在 $\mathscr{I}^{\pm}$ 上没有通量；而完整的 Boulware 真空在视界上奇异这一点，在中子星情形中不会困扰我们，因为那里根本不存在任何视界。

为了绝对确信我们为现实的黑洞正确选择了真空态，我们应当考虑在一个过去近乎 Minkowski 时空、未来为 Schwarzschild 时空的时空中发生的引力坍缩，如图 5.17 所示。如果真空在 $\mathscr{I}^{-}$ 上取标准的 Minkowski 形式，我们就可以追问这些模式如何穿过坍缩几何传播到 $\mathscr{I}^{+}$，这就定义了一个如式 (9.45) 中的 S 矩阵，用以确定渐近观者会看到什么。Hawking 最初发现黑洞辐射时所做的正是这件事；计算涉及一些烦琐的代数，但基本上是直截了当的，得到的温度与我们上面推导的相同。

当然，从完整的计算中我们能学到的不仅仅是黑体温度；我们会想问，比如说，当发射辐射的波长与 Schwarzschild 半径相当时会发生什么——在这种情形下我们的近似显然失效了。如果我们仔细研究任意种类的粒子从任何一种黑洞（也就是同时容许电荷和自旋）的发射，就会发现发射辐射的谱具有如下形式

\begin{equation*}
\langle \hat{n}_{\omega} \rangle = \frac{\Gamma(\omega)}{e^{2\pi(\omega - \mu)/\kappa} \pm 1}.
\tag{9.174}
\end{equation*}

这里的 $\kappa$ 当然就是表面引力。参数 $\mu$ 是一个化学势（chemical potential），刻画黑洞丢掉其守恒量子数的倾向；带电黑洞优先发射与黑洞电荷同号的粒子，而转动黑洞优先发射与黑洞角动量同号的粒子。因此，霍金辐射趋向于把黑洞带回 Schwarzschild 状态。$\Gamma(\omega)$ 是一个灰体因子（greybody factor），可以把它想象成波包被引力场反散射回黑洞而产生的。在高频极限下波长非常小，反散射可以忽略；在极低的频率下，波长变得大于 Schwarzschild 半径，反散射就变得重要起来。虽然灰体因子的解析表达式很难导出，但在频率很大和很小的极限情形下，标量场的灰体因子满足

\begin{align*}
\Gamma(\omega) &\to 1, \qquad \omega \gg \frac{1}{GM} \\
\Gamma(\omega) &\to \frac{A}{4\pi} \omega^{2}, \qquad \omega \ll \frac{1}{GM},
\tag{9.175}
\end{align*}

其中 $A$ 是黑洞的面积。

从经典广义相对论的观点来看，黑洞发出热辐射的发现当然是令人惊讶的，因为我们曾经强调过，从事件视界内部的点逃逸到无限远是不可能的。理解眼下情形的一种形象方式，是用 Feynman 图来思考真空涨落：涨落由不断地出现又消失的虚粒子/反粒子对来表示。这幅图像也很有帮助，例如有助于理解诸如 Lamb 位移（Lamb shift）这样的观测现象——原子光谱会受到光子与虚电子/正电子对相互作用的影响。通常情况下，这些粒子对总会湮灭，它们的效应只是间接的，通过与虚粒子相耦合的过程的重整化体现出来。然而，在存在事件视界时，虚粒子对中的一个成员偶尔会落入黑洞，而它的同伴则逃逸到无限远，如图 9.2 所描绘。在这幅图像中，我们观测为霍金辐射的正是这些逃逸的虚粒子。虚粒子对的总能量必须相加为零，但从无限远来看，落入的粒子可以带有负能量，因为渐近类时的 Killing 矢量在视界内部是类空的。这幅图像有些不正规，但它为正在发生的事情提供了一种有用的启发。

\begin{figure}[H]
\centering
\includegraphics[width=0.72\textwidth]{fig_9.2.pdf}
\caption*{图 9.2\quad 真空涨落偶尔会导致粒子/反粒子对中的一个落入事件视界，而另一个则作为霍金辐射逃逸到无限远。}
\end{figure}

一旦知道了黑洞温度的公式，我们就能够确定黑洞参数与热力学变量之间各关系式中的比例常数，如式 (6.118) 所列。霍金辐射可以说最终圆满成就了黑洞力学与热力学的联姻；稳态黑洞的行为就如同能量 $E = M$、温度 $T = \kappa/2\pi$、熵 $S = A/4G$ 的热平衡物体。这确实是极其巨大的熵。对宇宙中的物质场而言，熵近似等于相对论性粒子的数目；在一个 Hubble 半径之内，这个数目算出来是

\begin{equation*}
S_{\mathrm{M}} \sim 10^{88}.
\tag{9.176}
\end{equation*}

另一方面，黑洞的熵是以 Planck 单位量度的视界面积（记住，我们从头到尾都在取 $\hbar = 1$）。我们可以换算成天体物理单位，得到

\begin{equation*}
S_{\mathrm{BH}} \sim 10^{90} \left( \frac{M}{10^{6} M_{\odot}} \right)^{2}.
\tag{9.177}
\end{equation*}

这样，单个百万太阳质量的黑洞（例如在我们自己银河系的中心以及许多其他星系的中心就能找到）所具有的熵，就比可见宇宙中全部物质的熵还要多。宇宙的总熵远比我们本来能把它做到的要小——只要把更多的质量投入黑洞就行。（当宇宙学家说熵 $S_{\mathrm{M}}$ 很大时，他们的意思是：在一个曲率半径之内竟然存在这么多熵，这很令人惊讶。）我们之所以处在这样一个低熵状态，其原因想必与初始条件有关，或许还与暴胀有关。

回到黑洞力学，我们看到一个难题：宏观黑洞的熵会非常巨大，但从统计力学的观点来看，熵本应量度可及态数目的对数。经典黑洞由少数几个参数（质量、电荷和自旋）确定，因此很难知道那些态可能是什么。尽管如此，我们可以采取这样的态度：这个矛盾其实无关紧要，因为关于黑洞状态的任何信息想必都藏在事件视界背后。

把量子力学包括进来，使这个难题变得更糟而非更好，因为黑洞不仅会辐射，而且会蒸发。当我们开始研究弯曲时空中的 QFT 时，我们定下的规则之一就是：假定背景度规固定不变，不去操心量子场自身的能动张量会产生什么影响。尽管如此，即使在量子力学中我们也有能量守恒（例如，表现为渐近平直时空中守恒的 ADM 质量）。因此，当霍金辐射逃逸到无限远时，我们可以放心地断定，它会把能量从黑洞带走，黑洞的质量必然因此缩减。（这个现象并不违反面积定理，因为量子场的能动张量在视界附近将不满足弱能量条件。）随着质量缩减，表面引力增大，温度也随之升高；这是一个失控的过程，全部质量会在有限的时间内蒸发殆尽。代入数值，可以得到量级如下的寿命

\begin{equation*}
\tau_{\mathrm{BH}} \sim \left( \frac{M}{m_{\mathrm{P}}} \right)^{3} t_{\mathrm{P}} \sim \left( \frac{M}{M_{\odot}} \right)^{3} \times 10^{71}\,\mathrm{sec},
\tag{9.178}
\end{equation*}

其中 $m_{\mathrm{P}} \sim 10^{-5}\,\mathrm{g}$ 是 Planck 质量，$t_{\mathrm{P}} \sim 10^{-43}\,\mathrm{sec}$ 是 Planck 时间。由于 Hubble 时间为 $H_{0}^{-1} \sim 10^{18}\,\mathrm{sec}$，一个太阳质量黑洞的寿命约为宇宙年龄的 $10^{53}$ 倍。这看起来是很长的时间，但我们在这里谈论的是原理性的问题。

你可以看出为什么黑洞熵的问题变得如此严峻：一旦黑洞蒸发殆尽，我们就不能再求助于事件视界来隐藏黑洞那些所谓的态。黑洞不复存在，只剩下它产生的霍金辐射。这种辐射被认为应当是严格热性的（出射粒子中没有隐藏的关联），这一事实意味着，它没有任何办法传递巨量的信息，而指定我们的熵计算所隐含的那些态正需要这些信息。这样一来，如果我们把两个非常不同的初始态分别坍缩成两个质量、电荷和自旋都相同的黑洞，它们就会辐射成两团不可分辨的霍金粒子云。系统在变成黑洞之前、为规定其状态而投入的信息，看来已被抹去；这就是\textbf{信息丢失佯谬}（information loss paradox）。量子场论和广义相对论都以幺正演化为特征——指定早期时刻的一个态所需的信息，与指定较晚时刻的一个态所需的信息严格相等，因为二者通过运动方程联系在一起。但在把 QFT 与 GR 结合的过程中，这种幺正性显然被破坏了。看来很可能我们在论证的某个地方作了某个不恰当的假设，但很难看出在哪里。

传达信息丢失佯谬实质的一个途径，是考察图 9.3 所示的设想中的蒸发黑洞共形图。我们并不真正知道完整的时空应当是什么样子，但这里我们作了似乎合理的假设：形成一个奇点，以及一个与之相伴随的事件视界，二者在黑洞完全蒸发后都消失，留下一个具有 Minkowski 型因果结构的时空。如果我们用 Cauchy 面来思考，问题就显而易见了。一张从类空无限远 $i^{0}$ 延伸到奇点过去侧某个 $r = 0$ 点的非编时面，其未来依赖域将是整个时空，因此这样的面会是一个 Cauchy 面。但是一张延伸到奇点未来侧某个 $r = 0$ 点的类似的面就不是 Cauchy 面了，因为事件视界背后的区域不在它的依赖域之中。因此，由于信息消失于奇点，过去无法由未来反推。换句话说，这个过程看起来是时间不可逆的（在微观意义上，而不仅仅是统计意义上），尽管用来预言它的动力学定律在时间反演下是完全不变的。

\begin{figure}[H]
\centering
\includegraphics[width=0.72\textwidth]{fig_9.3.pdf}
\caption*{图 9.3\quad 设想中的蒸发黑洞共形图。能量被霍金辐射带走，黑洞最终完全蒸发殆尽，留下具有 Minkowski 时空因果结构的未来。落过事件视界进入奇点的信息看起来丢失了。}
\end{figure}

在讨论信息丢失佯谬时，请记住我们对黑洞蒸发的分析只是在量子场论与广义相对论相耦合的混合理论中进行的，并不是在现实的量子引力理论中。真实世界里究竟可能发生着什么呢？一种可能性是信息真的会丢失，幺正性真的会被破坏，而我们只得学会与之共处。许多物理学家觉得，引入这样一种可预见性的根本失效令人难以接受，而且也有论证表明，幺正性的破坏必然会导致能量守恒的破坏。另一种可能性是，幺正性在我们的世界中看似被破坏，只是因为进入黑洞的信息以某种方式逃逸到了空间一个不相连通的区域——一个婴儿宇宙（baby universe）。广义相对论预言黑洞中心存在奇点，而不是产生一个不相连通的区域；但显然我们正处于量子效应将剧烈改变经典预期的领域，所以我们应当保持开放的心态。

反对信息丢失的一些证据来自弦论。弦论自然地在 10 或 11 维时空中定义，其特征不仅在于一维的延展客体（弦），还在于统称为“膜”（branes）的各种更高维的延展客体。弦论的一个关键方面，是将玻色子与费米子联系起来的高度超对称性。在现实世界中，超对称性如果真的存在，就必须是自发破缺的，因为我们没有观测到与电子具有相同质量和电荷的玻色版本。但作为思想实验的工具，超对称性是无价的。人们可以构造出弦和膜的超对称位形，用来描述各种维度中的黑洞几何。弦论中有一个自由参数（实际上是一个标量场），即弦耦合（string coupling），它控制着引力的强度以及其他各种力的强度。如果我们考虑在某个弦耦合取值下描述黑洞的位形，那么当我们减小耦合时，Schwarzschild 半径最终会缩小到这个位形的尺寸以下，于是该位形就变成一团弱耦合的弦和膜。由于超对称性程度很高，我们可以确信，当我们改变弦耦合时，态的各种特征保持不变；特别是，我们预期自由度的数目（因而熵）不会改变。但在弱耦合区域没有黑洞，我们只有一团由通常自由度组成的“气体”（诚然，是高维中延展客体的气体），它的熵我们应该能够可靠地计算。

Strominger 和 Vafa 对一类特殊的五维超对称黑洞（带有不同种类的荷）考虑了这个过程。\footnote{A. Strominger and C. Vafa, “Microscopic origin of the Bekenstein-Hawking entropy,” \emph{Phys. Lett.} B \textbf{379}, 99 (1996), \texttt{http://arxiv.org/hep-th/9601029}。评述可参见 Johnson (2003)，或 A.W. Peet, “TASI lectures on black holes in string theory,” (2001), \texttt{http://arxiv.org/hep-th/0008241}。}他们得到了一个非凡的结果：系统在弱耦合下的自由度数目，与根据强耦合下的黑洞熵所预言的数目精确相符。由于黑洞熵非平庸地依赖于位形的荷，这种符合看来不大可能是偶然的。后续的研究把这一分析扩展到了其他类型的黑洞，并且不断地发现相符的结果。此外，通过考虑在弱耦合系统上的散射，我们甚至可以计算出预期的黑洞灰体因子；同样，结果与强耦合下的预期相符。因此，至少在弦论中，我们有极好的理由相信，黑洞辐射所隐含的自由度是真实存在的。

遗憾的是，弦论对态的计数，对于“黑洞状态的信息如何以某种方式传给出射的霍金辐射”这个问题，几乎没有提供直接的理解。尽管如此，我们当然应当认真对待事情果真如此的可能性，即使想象这样一个过程究竟如何实际运作还面临着严重的困难。当我们考虑把某些信息——或许以一套百科全书的形式——扔进一个大黑洞，而那时离黑洞蒸发殆尽还早得很的时候，困难就出现了。在这个阶段，黑洞温度很低，表面引力极小，事件视界附近的时空曲率也相当小。从百科全书的角度看，视界处没有任何特别的事情发生，我们应当预期它基本上毫发无损地落过视界。特别是，很难想象百科全书中的信息如何能够传递给早期发射出的霍金辐射。在幺正演化中，信息不可能被复制；它要么随百科全书一起落过视界，要么就必须在视界即将被穿越之前被有效地提取出来，而这似乎不合情理。我们或许可以指望，信息伴随着百科全书进入奇点附近的某个区域，并以某种方式保存在那里，直到黑洞变得很小的晚期。但到那时，大部分辐射粒子已经被发射出去了，最后一阵辐射所能到达的态的数目，一般会小于描述可能已经落入黑洞的各种不同态所需要的数目。

因此，要想认为信息以某种方式编码在出射辐射之中，看来就有必要在很早的时候就把关联编码到霍金粒子之中。我们刚刚论证过这一点很难做到，因为黑洞很大时，视界是一个平淡无奇的地方。摆脱这一困境的一个可以设想的办法，是迈出戏剧性的一步：放弃局域量子场论。换句话说，我们一直隐含地假设，信息可以被合理地描述为定域在空间的某个区域之中；这是通常量子场论的一个无可争议的特征。但量子引力也许有所不同，黑洞所包含的信息或许以某种方式非局域地弥散在视界之上。这个建议本身并不能直接给出把信息送进出射霍金辐射的机制，但它确实使我们给出的、关于为什么这样做会很困难的一些论证受到质疑。

非局域性的一种具体实现，以\textbf{全息原理}（holographic principle）为名。这个思想最初由 't Hooft 和 Susskind 提出：空间一个区域中的自由度数目，不是与该区域的体积成正比（正如在局域场论中所预期的那样），而是与该区域边界的面积成正比。\footnote{评述可参见 R. Bousso, “The Holographic Principle” (2002), \texttt{http://arxiv.org/hep-th/0203101}。}其灵感当然来自黑洞熵，后者按事件视界的面积标度；如果熵计数的是可及态的数目，全息就能解释为什么起作用的是面积，而不是所包围的体积。你也许会担心如何处理闭合宇宙的情形：在其中某个区域可能几乎由整个空间构成，边界却非常小；不过，只要把空间区域换成一组从边界向内延伸的“光片”（lightsheets），就可以表述出一个更具协变性的全息原理版本。全息原理最伟大的胜利是在第 8 章提到过的 AdS/CFT 对应中。在那里，反 de Sitter 背景上的量子引力物理，等价于定义在 AdS 边界上、维数低一维的无引力共形场论。可以设想，我们在宇宙中观测到的所有物理现象，都可以由某个定义在更低维数中的普通非引力理论的非局域全息投影来描述；我们究竟应当如何构造这样一种对应、又如何将它与观测联系起来，目前还完全不清楚，但宇宙学以及宇宙大尺度结构的考虑也许会是一个有希望的切入点。

关于黑洞熵、弦论和全息原理的这些评论，显然无意作为对一个非常活跃的研究领域的细致导引。毋宁说，它们意在指出引力物理前沿正在探索的若干可能性。经典广义相对论是迄今发明的最优美的物理理论，但我们完全有权利期待，GR 与物理学其他领域的综合，将揭示出我们如今只能想象的层层之美。
